\documentclass[a4paper,12pt]{amsart}
\usepackage{hyperref}
\usepackage{amsmath,amsthm,amssymb,amscd}
\usepackage{color}

\title[Lipschitz homotopy convergence II]{Lipschitz homotopy convergence of Alexandrov spaces II}

\author[T. Fujioka]{Tadashi Fujioka}
\address[T. Fujioka]{Department of Mathematics, Osaka University, Osaka 560-0043, Japan}
\email{fujioka@cr.math.sci.osaka-u.ac.jp}

\author[A. Mitsuishi]{Ayato Mitsuishi}
\address[A. Mitsuishi]{Department of Applied Mathematics, Fukuoka University, Fukuoka 814-0180, Japan}
\email{mitsuishi@fukuoka-u.ac.jp}

\author[T. Yamaguchi]{Takao Yamaguchi}
\address[T. Yamaguchi]{Department of Mathematics, University of Tsukuba, Tsukuba 305-857, Japan}
\email{takao@math.tsukuba.ac.jp}

\date{\today}
\subjclass[2020]{53C20, 53C23}
\keywords{Alexandrov spaces, Gromov-Hausdorff convergence, noncollapsing, Lipschitz homotopy, gradient flows, good coverings, extremal subsets}

\theoremstyle{plain}
\newtheorem{thm}{Theorem}
\newtheorem{cor}[thm]{Corollary}
\newtheorem{prop}[thm]{Proposition}
\newtheorem{lem}[thm]{Lemma}
\newtheorem{clm}[thm]{Claim}
\newtheorem{subclm}[thm]{Subclaim}
\newtheorem{comp}[thm]{Complement}

\theoremstyle{definition}
\newtheorem{dfn}[thm]{Definition}
\newtheorem{rem}[thm]{Remark}
\newtheorem{prob}[thm]{Problem}
\newtheorem{conv}[thm]{Convention}

\numberwithin{thm}{section}
\numberwithin{equation}{section}

\DeclareMathOperator{\dist}{dist}
\DeclareMathOperator{\diam}{diam}
\DeclareMathOperator{\vol}{vol}
\DeclareMathOperator{\supp}{supp}

\newcommand{\step}[1]{\medskip\noindent\textit{#1.}}

\begin{document}

\begin{abstract}
We obtain a quantitative version of the Lipschitz homotopy convergence introduced in [Mitsuishi and Yamaguchi 2019b] for a moduli space of compact Alexandrov spaces without collapsing. 
We also show that the Lipschitz homotopies can be chosen to preserve extremal subsets.
\end{abstract}

\maketitle

\section{Introduction}

\subsection{Main result}

This paper is a continuation of \cite{MY:lip} by the second and third authors.
The main result of \cite{MY:lip} was that the class of Alexandrov spaces with upper bounds on dimension and diameter and with lower bounds on curvature and volume contains only finitely many Lipschitz homotopy types.
It is natural to expect that there exists a uniform Lipschitz constant for the Lipschitz homotopies, and we will solve this problem in this paper.

Let $X$ and $Y$ be metric spaces.
For a constant $C>0$, we define the {\it $C$-Lipschitz homotopy distance} $d_{C\mathchar`-\mathrm{LH}}(X,Y)$ between $X$ and $Y$ as the infimum of $\epsilon>0$ such that there are $\epsilon$-approximations $f:X\to Y$ and $g:Y\to X$ which are $C$-Lipschitz, together with $C$-Lipschitz maps $F:X\times [0,\epsilon]\to X$ and $G:Y\times [0,\epsilon]\to Y$ satisfying
\[F(\cdot,0)=g\circ f,\quad F(\cdot,\epsilon)=1_X,\quad G(\cdot,0)=f\circ g,\quad G(\cdot,\epsilon)=1_Y,\]
where $1$ denote the identity maps.
We refer to such $f$ and $g$ as \textit{$(C,\epsilon)$-Lipschitz homotopy approximations} between $X$ and $Y$.
In other words, they are $\epsilon$-approximations that are also $C$-Lipschitz homotopy equivalence whose time interval has length $\epsilon$.

For $n\in\mathbb N$ and $D,v>0$, let $\mathcal A(n,D,v)$ denote the class of compact $n$-dimensional Alexandrov spaces with curvature $\ge-1$, diameter $\le D$, and volume $\ge v$.
Recall that $\mathcal A(n,D,v)$ is compact with respect to the usual Gromov-Hausdorff distance, denoted by $d_\mathrm{GH}$.

The following is the main result of the present paper.

\begin{thm}\label{thm:main}
For given $n$, $D$, and $v$, there exists a constant $C\gg 1$ satisfying the following:
for any $M,M'\in\mathcal A(n,D,v)$ with Gromov-Hausdorff distance less than $C^{-1}$, we have
\[d_{C\mathchar`-\mathrm{LH}}(M,M')<Cd_\mathrm{GH}(M,M').\]
More precisely, for any given $\epsilon$-approximation between $M$ and $M'$, where $\epsilon<C^{-1}$, there exists a $(C,C\epsilon)$-Lipschitz homotopy approximation between $M$ and $M'$ that is $C\epsilon$-close to it.
\end{thm}

As a direct consequence, we have 

\begin{cor}\label{cor:main}
$\mathcal A(n,D,v)$ is compact with respect to the $C$-Lipschitz homotopy distance.
It contains only finitely many $C$-Lipschitz homotopy types.
\end{cor}

It should be emphasized that our results are new even in the Riemannian case.
Theorem \ref{thm:main} and Corollary \ref{cor:main} are closely related to Perelman's stability theorem (\cite{Per:alex}, \cite{K:stab}), which asserts that under the same assumption as in Theorem \ref{thm:main}, $M$ and $M'$ are homeomorphic.
In particular, it implies the finiteness of homeomorphism types in $\mathcal A(n,D,v)$.
Although Perelman claimed that the stability homeomorphism can be chosen to be bi-Lipschitz, the proof has never been published.

\subsection{Strategies}

There are two ways of proving the main theorem.
One is the same strategy as in the previous paper \cite{MY:lip}, which uses the good coverings of Alexandrov spaces introduced in \cite{MY:good}.
The other is a Lipschitz analogue of Petersen's strategy in \cite{Peter}, which relies only on local contractibility.
Both have their own advantages (and disadvantages), so we will provide two proofs of the main theorem.

In both cases, the key geometric property is the uniform local Lipschitz contractibility of Alexandrov spaces, Theorem \ref{thm:conv}.
It asserts that any $\epsilon$-ball in an Alexandrov space in $\mathcal A(n,D,v)$ admits a convex hull with a $C$-Lipschitz contraction of time $\epsilon$, where $C$ depends only on $n$, $D$, and $v$.
This is a strengthening of an earlier result of Perelman and Petrunin \cite[4.3]{PP:ext} on convex hulls in Alexandrov spaces (cf.\ \cite[7.1.3]{Pet:sc}).
The proof is a combination of two techniques: the construction of a strictly concave function by Perelman and Petrunin \cite[4.3]{PP:ext} and the stability of Lipschitz contractible balls by the second and third authors \cite[1.2]{MY:stab}.

In the first strategy using good coverings, we need an additional geometric argument based on Theorem \ref{thm:conv} to construct a quantitative version of a good covering (Theorem \ref{thm:good}).
Once the existence of such a good covering is shown, the rest of the proof is almost the same as in the previous paper \cite{MY:lip}.
We prove that any metric space with a quantitative good covering is uniformly Lipschitz homotopy equivalent to a geometric realization of the nerve (Theorem \ref{thm:nerve}).

In the second strategy based only on local contractibility, we prove the Lipschitz version of Petersen's stability theorem \cite[Theorem A]{Peter} (Theorem \ref{thm:Peter}).
The proof is along the same line as the original one.
However, the existence of a Lipschitz domination of an Alexandrov space by a polyhedron needs an additional argument not found in \cite{Peter} (Proposition \ref{prop:dom}).
Fortunately this is done by modifying an argument in the previous paper \cite{MY:lip}.

\subsection{Modifications}

There are two modifications of the main theorem.
One is the gluing with an almost isometry as in \cite[1.4]{MY:lip}.
Recall that there is an almost isometry between closed domains of strained parts of two close Alexandrov spaces (Theorem \ref{thm:reg}; see \cite[\S2.3]{MY:lip} or \cite{BGP} for the definition and properties of strainers).

For an $n$-dimensional Alexandrov space $M$, let $\mathcal R_M(\delta,\ell)$ denote the set of points in $M$ with $(n,\delta)$-strainers of length $>\ell$, where $\ell\le 1$ and $\delta$ is less than a sufficiently small constant depending only on $n$.
We also denote by $\tau(\delta)$ a positive function depending only on $n$ tending to $0$ as $\delta\to0$.

\begin{thm}\label{thm:isom}
Under the same assumption as in Theorem \ref{thm:main}, the Lipschitz homotopy approximation from $M$ to $M'$ can be chosen to be a $\tau(\delta)$-almost isometric open embedding on $\mathcal R_M(\delta,\ell)$, provided that the Gromov-Hausdorff distance is less than $C^{-1}\delta\ell$.
\end{thm}

The proof is along the same line as in the previous paper \cite{MY:lip}, but we need to check the Lipschitz constant more carefully.

The other modification is new, which is the relative version of the main theorem with respect to extremal subsets.
Extremal subsets are singular sets of Alexandrov spaces introduced by Perelman and Petrunin \cite{PP:ext}.
It is known that in some sense any Alexandrov space is uniquely stratified into extremal subsets.
Kapovitch \cite[\S9]{K:stab} proved that the stability homeomorphism of Perelman can be chosen to preserve such extremal subsets.
It is natural to expect that the same is true for our Lipschitz homotopy.

As before, let $f:X\to Y$ and $g:Y\to X$ be $(C,\epsilon)$-Lipschitz homotopy approximations with $C$-Lipschitz homotopies $F:X\times[0,\epsilon]\to X$ and $G:Y\times[0,\epsilon]\to Y$.
For $A\subset X$ and $B\subset Y$, we say that $f$ and $g$ \textit{respect} $A$ and $B$ if
\[f(A)\subset B,\quad g(B)\subset A,\quad F(A,t)\subset A,\quad G(B,t)\subset B\]
for all $0\le t\le\epsilon$.

\begin{thm} \label{thm:ext}
Under the same assumption as in Theorem \ref{thm:main}, let $E$ and $E'$ be extremal subsets of $M$ and $M'$, respectively, that are $C^{-1}$-close to each other.
Then the Lipschitz homotopy approximations between $M$ and $M'$ can be chosen to respect $E$ and $E'$.
\end{thm}

Since a limit of extremal subsets is also extremal, the above situation naturally occurs under a convergence of Alexandrov spaces.
Note that this theorem is independent of the previous one because any proper extremal subset is contained in a non-strained part of the ambient Alexandrov space.

We will give two different proofs of Theorem \ref{thm:ext}.
In the first proof the key fact is that any extremal subset is invariant under the gradient flows of semiconcave functions.
Since the Lipschitz homotopies of Theorem \ref{thm:main} are given by such gradient flows, it turns out that they automatically preserve extremal subsets.
The second proof is more general, which provides a way of deforming any given Lipschitz homotopy approximation to preserve an extremal subset.
This uses a neighborhood deformation retraction to an extremal subset constructed by the first author \cite[1.5]{F:reg}.

\step{Organization}
Section \ref{sec:pre} contains background materials.
Section \ref{sec:conv} is devoted to the proof of the key theorem \ref{thm:conv}.
In Section \ref{sec:good}, using it, we construct a quantitative version of a good covering (Theorem \ref{thm:good}).
These two sections are the main geometric part of the present paper.
Section \ref{sec:main1} provides the first proof of Theorem \ref{thm:main} using good coverings, which is almost parallel to \cite[\S3]{MY:lip}.
Section \ref{sec:main2} provides the second proof of Theorem \ref{thm:main} using only local contractibility, which is almost parallel to \cite[\S2--\S4]{Peter}, except Proposition \ref{prop:dom}.
In Section \ref{sec:isom}, we prove Theorem \ref{thm:isom}.
Finally in Section \ref{sec:ext}, we give two proofs of Theorem \ref{thm:ext}.

\step{Acknowledgments}
The first author was supported by JSPS KAKENHI Grant Number 22J00100.
The second author was supported by JSPS KAKENHI Grant Number 20K03598.
The third author was supported by JSPS KAKENHI Grant Number 21H00977.


\section{Preliminaries} \label{sec:pre}

In this section we fix notation, introduce terminology on Lipschitz homotopy and Gromov-Hausdorff convergence, and review the properties of gradient flows.
We refer to \cite{BGP} for Alexandrov spaces.

\subsection{Notation}

The distance between $x$ and $y$ will be denoted by $|xy|$.
The open $r$-ball and $r$-sphere around a point $p$ will be denoted by $B(p,r)$ and $S(p,r)$, respectively.
Similarly, the open $r$-neighborhood and its boundary of a subset $A$ will be denoted by $B(A,r)$ and $S(A,r)$.
For a metric space $X$, the same space with metric multiplied by $\lambda$ is denoted by $\lambda X$.
We denote by $1_X$ the identity of $X$.

The distance function to a point $p$ will be denoted by $|p\cdot|$ or $\dist_p$.
The direction of a shortest path from $p$ to $q$ is denoted by $\uparrow_p^q$ and the set of directions of all shortest paths from $p$ to $q$ is denoted by $\Uparrow_p^q$.
We use the same notation for a subset $A$, i.e., $|A\cdot|$, $\dist_A$, $\uparrow_p^A$, and $\Uparrow_p^A$.
The angle and comparison angle will be denoted by $\angle$ and $\tilde\angle$, respectively.
The space of directions and tangent cone at $p$ will be denoted by $\Sigma_p$ and $T_p$, respectively.

Throughout the paper, $n$ denotes the dimension, $D$ denotes an upper diameter bound, and $v$ denotes a lower volume bound.
The lower curvature bound is always $-1$.
We use the symbol $C$ to denote various large positive constants, usually depending only on $n$, $D$, and $v$.
Another symbol $\epsilon$ is also used to denote a small positive number, which is usually assumed to be less than some small constant depending only on $n$, $D$, and $v$.
The usage of these symbols will be explained in more detail in Sections \ref{sec:main1} and \ref{sec:main2} (see Conventions \ref{conv:C} and \ref{conv:C'}).

\subsection{Lipschitz homotopy}

We introduce the quantitative version of Lipschitz homotopy used in \cite{MY:lip}.
Fix $C>0$ and $\epsilon>0$ (usually $C\gg1$ and $\epsilon\ll1$, but not necessarily assumed here).
Loosely speaking, we will add the word ``$(C,\epsilon)$-Lipschitz'' to the terminology of homotopy theory if the homotopy under consideration is $C$-Lipschitz continuous and its time interval has length $\epsilon$.

Let $X$ and $Y$ be metric spaces.
For simplicity, whenever we consider a product space with a time interval, such as $X\times[0,\epsilon]$, we equip it with the $L_1$ distance.

\begin{dfn}
Two maps $h_0, h_\epsilon:X\to Y$ are said to be \textit{$(C,\epsilon)$-Lipschitz homotopic} if there exists a $C$-Lipschitz map $h:X\times[0,\epsilon]\to Y$ such that $h_i = h(\cdot,i)$ for $i=0,\epsilon$.
The map $h$ is called a \textit{$(C,\epsilon)$-Lipschitz homotopy}.

We say that $X$ and $Y$ are \textit{$(C,\epsilon)$-Lipschitz homotopy equivalent} if there are $C$-Lipschitz maps $f:X\to Y$ and $g:Y\to X$ such that $g\circ f$ and $f\circ g$ are $(C,\epsilon)$-Lipschitz homotopic to $1_X$ and $1_Y$, respectively.
The maps $f$ and $g$ are called \textit{$(C,\epsilon)$-Lipschitz homotopy equivalences}.
\end{dfn}

Note that if $h$ is a $(C,C'\epsilon)$-Lipschitz homotopy, where $C'\ge 1$, then $h'(\cdot,t):=h(\cdot,C't)$, where $t\in[0,\epsilon]$, is a $(CC',\epsilon)$-Lipschitz homotopy.
Hence any large constant multiple of the time (independent of $\epsilon$) can be absorbed in the Lipschitz constant.
This basic fact will be used  implicitly throughout the paper.

The following is a special case of Lipschitz homotopy equivalence.

\begin{dfn}
We say that $A\subset X$ is a \textit{$(C,\epsilon)$-Lipschitz strong deformation retract} of $X$ if there is a $C$-Lipschitz map $F:X\times[0,\epsilon]\to X$ such that $F(x,0)=x$, $F(x,\epsilon)\in A$, and $F(a,t)=a$ for any $x \in X$, $a \in A$, and $0\le t\le \epsilon$. 
The map $F$ is called a \textit{$(C,\epsilon)$-Lipschitz strong deformation retraction} of $X$ to $A$.

Moreover, if $A$ is a point, we say that $X$ is \textit{$(C,\epsilon)$-Lipschitz strongly contractible} and $F$ is a \textit{$(C,\epsilon)$-Lipschitz strong contraction}.
In this case $A$ is called the \textit{center} of the contraction.
\end{dfn}

We will often omit the word ``strong'': in fact we will not deal with ``weak'' ones (i.e., not fixing $A$) in this paper.
Note that if $X$ is $(C,\epsilon)$-Lipschitz contractible, then the diameter of $X$ is at most $2C\epsilon$.

We also use the following relative version.
Let $A\subset X$ and $B\subset Y$.

\begin{dfn}
Let $h:X \times [0,\epsilon] \to Y$ be a $(C,\epsilon)$-Lipschitz homotopy.
We say that $h$ \textit{respects} $A$ and $B$ if $h(A,t)\subset B$ for any $0\le t\le\epsilon$.

We say that the pairs $(X,A)$ and $(Y,B)$ are {\it $(C,\epsilon)$-Lipschitz homotopy equivalent} if there are $(C,\epsilon)$-Lipschitz homotopy equivalences $f:X\to Y$ and $g:Y\to X$ with $f(A)\subset B$ and $g(B)\subset A$ whose homotopies respecting $A$ and $B$.
In this case, $f$ and $g$ are said to \textit{respect} $A$ and $B$.
\end{dfn}

\subsection{Gromov-Hausdorff convergence}
We recall the terminology of Gromov-Hausdorff convergence.

A map $f:X\to Y$ is called an \textit{$\epsilon$-approximation} if it satisfies
\begin{enumerate}
\item $||f(x)f(x')| - |xx'||<\epsilon$ for all $x,x'\in X$; 
\item for any $y\in Y$, there is $x\in X$ such that $|f(x)y|<\epsilon$.
\end{enumerate}

If $f:X\to Y$ is an $\epsilon$-approximation, we say that $f(p)\in Y$ is a \textit{lift} of $p\in X$.
Similarly, we say that $\dist_{f(p)}$ is a \textit{lift} of $\dist_p$.
In this way, any objects on $X$ defined by distance functions can be lifted to $Y$ by replacing $\dist_p$ with $\dist_{f(p)}$.

The \textit{Gromov-Hausdorff distance} $d_\mathrm{GH}(X,Y)$ between $X$ and $Y$ is defined as the infimum of $\epsilon>0$ such that there are $\epsilon$-approximations $f:X\to Y$ and $g:Y\to X$.
Whenever we discuss the Gromov-Hausdorff distance, we will implicitly fix some approximation.

For $A\subset X$ and $B\subset Y$, if an $\epsilon$-approximation $f:X\to Y$ satisfies $f(A)\subset B(B,\epsilon)$ and $B\subset B(f(A),\epsilon)$, then we call $f$ an \textit{$\epsilon$-approximation} from $(X,A)$ to $(Y,B)$ (where $B(\cdot,\cdot)$ denote metric neighborhoods).
This defines the convergence of subsets under the Gromov-Hausdorff convergence of spaces.

If $(C,\epsilon)$-Lipschitz homotopy equivalences are also $\epsilon$-approximations, we call them \textit{$(C,\epsilon)$-Lipschitz homotopy approximations}.
This leads to the definition of the \textit{$C$-Lipschitz homotopy distance}, as discussed in the introduction.

\subsection{Gradient flows}

We briefly review the properties of the gradient flows of semiconcave functions in Alexandrov spaces.
See \cite{Pet:sc} for more details.
In what follows, for simplicity we assume that functions are defined on the whole space, but this is not necessary.

Let $M$ be an Alexandrov space and $f:M\to\mathbb R$.
If $M$ has no boundary, $f$ is said to be \textit{$\lambda$-concave}, where $\lambda\in\mathbb R$, if its restriction to any unit-speed shortest path $\gamma(t)$ is $\lambda$-concave, i.e., $f\circ\gamma(t)-(\lambda/2)t^2$ is concave.
If $M$ has boundary, we require the same condition for the tautological extension of $f$ to the double of $M$.
A \textit{semiconcave} function is locally a $\lambda$-concave function, where $\lambda$ depends on each point.
A \textit{strictly concave} function is a $\lambda$-concave function with $\lambda<0$.

A typical example of a semiconcave function is a distance function $\dist_A$, where $A\subset M$.
Note that the concavity of $\dist_A$ at $x\in M\setminus A$ is bounded above by a constant depending only on a lower bound of $|Ax|$ (and the lower curvature bound $-1$; see \cite[2.4]{MY:loc} for instance).

Let $f:M\to\mathbb R$ be a semiconcave function and $p\in M$.
The directional derivative $df=d_pf$ is a $0$-concave function defined on the tangent cone $T_p$ at $p$.
The \textit{gradient} $\nabla_pf\in T_p$ of $f$ at $p$ is characterized by the following two properties:
\begin{enumerate}
\item $df(v)\le\langle\nabla_pf,v\rangle$ for any $v\in T_p$;
\item $df(\nabla_pf)=|\nabla_pf|^2$.
\end{enumerate}
Here $|\cdot|$ denotes the norm, i.e., the distance to the vertex of the cone and $\langle\cdot,\cdot\rangle$ is the scalar product defined by the norm and the angle at the vertex as in Euclidean plane.
The concavity of $df$ guarantees the existence and uniqueness of the gradient.

The \textit{gradient curve} of a semiconcave function $f$ is a curve whose right tangent vector is unique and equal to the gradient of $f$.
For any $\lambda$-concave function $f$, the gradient curve $\Phi(x,t)$ starting at $x\in M$ exists and is unique for all $t\ge0$. 
The map $\Phi:M\times[0,\infty)\to M$ is called the \textit{gradient flow} of $f$.

The most important property of the gradient flow for this paper is its Lipschitz continuity.
Recall that we are considering the $L_1$ distance on a product space with a time interval.

\begin{prop}[{\cite[2.1.4]{Pet:sc}}]\label{prop:lip}
Let $f$ be a $\lambda$-concave function on $M$ and $\Phi:M\times[0,\infty)\to M$ its gradient flow.
Then $\Phi_t:=\Phi(\cdot,t):M\to M$ is $e^{\lambda t}$-Lipschitz for any $t\ge 0$.
In particular, if $f$ is $C$-Lipschitz, the restriction of $\Phi$ to $M\times[0,T]$ is $\max\{C,e^{\lambda T},1\}$-Lipschitz.
\end{prop}

We conclude this section with two basic lemmas regarding gradient which will be used throughout the paper.
The first one is a gradient estimate for a strictly concave function.

\begin{lem}\label{lem:grad}
Let $f$ be a $(-\lambda)$-concave function on $M$, where $\lambda>0$.
If $p\in M$ is the unique maximum point of $f$, then for any $x\in M\setminus\{p\}$, we have
\[|\nabla_xf|\ge df(\uparrow_x^p)\ge\lambda|px|/2.\]
\end{lem}

\begin{proof}
This immediately follows from the $(-\lambda)$-concavity:
\[df(\uparrow_x^p)\ge\frac{f(p)-f(x)+\lambda|px|^2/2}{|px|}.\]
\end{proof}

The second one is about the following regularity introduced in \cite{MY:stab}.
For $c>0$ and $A\subset M$, we say that $f$ is \textit{$c$-regular in the direction to $A$} if $df(\uparrow_x^A)>c$ for any $x\in M\setminus A$ and some shortest path from $x$ to $A$.

\begin{lem}[{\cite[2.3]{MY:stab}}]\label{lem:reg}
Let $f$ be a semiconcave function on $M$ that is $c$-regular in the direction to $A\subset M$.
Then the gradient flow of $f$ outside $A$ monotonically decreases the distance to $A$ with velocity at least $c$.
\end{lem}

\begin{proof}
This follows from the first variation formula and the property (1) of the gradient.
Indeed, if $\alpha(t)$ denotes a gradient curve of $f$, we have
\[\frac d{dt}|A\alpha(t)|=-\langle\nabla_{\alpha(t)}f,\Uparrow_{\alpha(t)}^A\rangle\le-df(\uparrow_{\alpha(t)}^A)<-c.\]
\end{proof}

\section{Lipschitz contractible convex hulls}\label{sec:conv}

In this section we prove the following key theorem.
Fix $n$, $D$, and $v$.

\begin{thm}[cf.\ {\cite[4.3]{PP:ext}}]\label{thm:conv}
There is $C=C(n,D,v)>0$ satisfying the following:
for any $M \in\mathcal A(n,D,v)$, $p \in M$, and $0<\epsilon<C^{-1}$, there exists a $(C,\epsilon)$-Lipschitz strongly contractible convex domain $U$ containing $B(p,\epsilon)$.

More precisely, the following hold:
\begin{enumerate}
\item there exists a $1$-Lipschitz $(-(C\epsilon)^{-1})$-concave function $h$ defined on a neighborhood of $p$ such that for any $0\le t\le C^{-1}$,
\[U(t):=\{h>-t\epsilon\}\]
is the desired domain;
\item $U(t)$ also contains the $(C^{-1}\epsilon)$-ball around its peak $q$, i.e., the unique maximum point of $h$;
\item for any $M'\in\mathcal A(n,D,v)$ that is $(C^{-1}\epsilon)$-close to $M$, there is a lift $h'$ of $h$ to $M'$ that satisfies the same properties as (1) and (2) for any lift $p'\in M'$ of $p$.
\end{enumerate}
\end{thm}

\begin{rem}
The center of the contraction is not necessarily $p$ or $q$.
\end{rem}

\begin{rem}
The contractibility implies that the diameter of $U$ is at most $2C\epsilon$.
The existence of such a convex hull of uniformly bounded diameter was proved by Perelman and Petrunin \cite[4.3]{PP:ext} (cf.\ \cite[7.1.3]{Pet:sc}).
Our result strengthens it.
\end{rem}

The detailed statement in the second half will be used in the next section to construct a quantitative version of a good covering, which will be used in the first proof of Theorem \ref{thm:main} in Section \ref{sec:main1}.
For the second proof in Section \ref{sec:main2}, the first half is enough.

The proof is a combination of techniques from \cite[4.3]{PP:ext} and \cite[1.2]{MY:stab}.
The construction of the strictly concave function $h$ is due to \cite[4.3]{PP:ext}.
We use the gradient flow of $h$ to contract $U$ to a tinier ball, while we need the stability of contractible balls as in \cite[1.2]{MY:stab} to crush this tinier ball.
Thus the first half of the proof is the same as that of \cite[4.3]{PP:ext} and the second half is the same as that of \cite[1.2]{MY:stab}.
For the convenience of the reader we include the complete proof.

\begin{proof}[Proof of Theorem \ref{thm:conv}]
We prove it by contradiction.
If the statement does not hold, there are $C_i\to\infty$, $0<\epsilon_i<C_i^{-1}$, $M_i\in\mathcal A(n,D,v)$, and $p_i\in M_i$ such that there does not exist such a strictly concave function $h_i$ satisfying the conclusion.
Instead of $M_i$, we will consider $\tilde M_i:=\epsilon_i^{-1}M_i$, and will write metric balls and spheres in $\tilde M_i$ as $\tilde B(\cdot,\cdot)$ and $\tilde S(\cdot,\cdot)$, respectively.
We will also consider $\tilde h_i:=\epsilon_i^{-1}h_i$ on $\tilde M_i$, for which the conditions (1) and (2) reduce to:
\begin{enumerate}
\item $\tilde h_i$ is a $1$-Lipschitz $(-C_i^{-1})$-concave function such that for any $0\le t\le C_i^{-1}$, the superlevel set $\tilde U_i(t):=\{\tilde h_i>-t\}$ is a $(C_i,1)$-Lipschitz contractible domain containing $\tilde B(p_i,1)$;
\item $\tilde U_i(t)$ also contains the $C_i^{-1}$-ball around its peak $q_i$.
\end{enumerate}
To obtain a contradiction, it suffices to show these conditions with $C_i$ replaced by a constant $C$ for some subsequence (see below for (3)).

Since $\epsilon_i\to 0$, we may assume that $(\tilde M_i,p_i)$ converges to a nonnegatively curved Alexandrov space $(N,p)$ of dimension $n$ in the pointed Gromov-Hausdorff topology.
Moreover, the existence of a uniform lower volume bound for $M_i$ implies that the ideal boundary $N(\infty)$ of $N$ has dimension $n-1$.
Indeed, by the Bishop-Gromov inequality, we see that the volume ratio $\vol_n(\tilde B(p_i,R))/R^n$ is uniformly bounded below by a constant depending only on $n$, $D$, and $v$ for any $i\gg R\gg 1$, where $\vol_n$ denotes the $n$-dimensional Hausdorff measure.

In what follows, we first prove the corresponding statement in $N$ and then show that the construction can be lifted to $\tilde M_i$.
Note that for any $M_i'\in\mathcal A(n,D,v)$ that is $(C_i^{-1}\epsilon_i)$-close to $M_i$, the rescaled sequence $\tilde M_i':=\epsilon_i^{-1}M_i'$ also converges to $N$.
Therefore the condition (3) obviously follows from the following proof.

\step{Step 1}
We first construct a strictly concave function $h$ defined on an arbitrarily large neighborhood of $p$, which can be lifted to a function $\tilde h_i$ on $\tilde M_i$ with the same concavity.
We show that a superlevel set of $h$ is a convex domain containing the unit ball around $p$ and the same holds for $\tilde h_i$.
This part is exactly the same as the proof of \cite[4.3]{PP:ext}.

Take a maximal $\delta$-discrete set $\{l_\alpha\}_{\alpha\in A}$ in $N(\infty)$, where $\delta>0$ will be determined in Claim \ref{clm:conc}.
For each $l_\alpha$ there exists a corresponding ray $r_\alpha$ emanating from $p$, which yields the Busemann function $b_\alpha$, i.e.,
\[b_\alpha:=\lim_{t\to\infty}|r_\alpha(t)\cdot|-t.\]

Choose $R\gg1$ large enough so that $R^{-1}S(p,R)$ is sufficiently close to $N(\infty)$ in the Gromov-Hausdorff distance.
Fix $x\in S(p,R)$.
Then there exists $A_x\subset A$, whose corresponding rays are near a shortest path $px$, such that
\begin{equation}\label{eq:a}
\# A_x\ge c/\delta^{n-1}
\end{equation}
for some $c>0$ independent of $x$, and
\begin{equation}\label{eq:b}
b_\alpha(p)\ge b_\alpha(x)+R/2
\end{equation}
for any $\alpha\in A_x$.
The former follows from the fact that $\dim N(\infty)=n-1$ and the Bishop-Gromov inequality.
Indeed, $c$ depends only on $n$, $D$, and $v$.

Next find $\hat R\gg R$ such that for any $\alpha\in A$ the normalized distance function $|q_\alpha\cdot|-\hat R$ from $q_\alpha:=r_\alpha(\hat R)$ is sufficiently close to $b_\alpha$.
In particular for any $\alpha\in A_x$ an inequality similar to \eqref{eq:b} holds, i.e.,
\[|q_\alpha p|\ge|q_\alpha x|+R/3.\]
We define a function $h_x$ on $B(p,2R)$ by
\[h_x:=\frac1{\# A_x}\sum_{\alpha\in A_x}\phi(|q_\alpha\cdot|),\]
where $\phi$ is a $C^2$ function such that $1/2\le\phi'\le 1$ and $\phi''=-(8R)^{-1}$ on $[\hat R-2R,\hat R+2R]$.

Clearly $h_x$ is $1$-Lipschitz, $h_x(p)\ge h_x(x)+R/6$, and $h_x(p)=h_{x'}(p)$ for $x\neq x'\in S(p,R)$.
Moreover,

\begin{clm}[{\cite[3.6]{Per:mor}}]\label{clm:conc}
$h_x$ is $(-\lambda)$-concave on $B(p,2R)$ for some $\lambda>0$ depending only on $c$ and $R$, provided $\delta$ is small enough.
\end{clm}

\begin{proof}
If $f$ is a $C^2$ function, then $(\phi\circ f)''=\phi''\circ f\cdot(f')^2+\phi'\circ f\cdot f''$.
Hence if $|f(0)-\hat R|<2R$, $|f'(0)|>c>0$, and $f''(0)<c^2/16R$, then $(\phi\circ f)''(0)<-c^2/16R$.
This estimate holds even if $f$ is not $C^2$.

For any direction $\xi$ at any point of $B(p,2R)$, a volume comparison argument using \eqref{eq:a} shows that  for the majority of $\alpha$, $|d\dist_{q_\alpha}(\xi)|$ is bounded from below by a constant depending only on $c$, provided $\delta$ is small enough.
Moreover, since $\hat R\gg R$ and $N$ has nonnegative curvature, $\dist_{q_\alpha}$ is almost $0$-concave on $B(p,2R)$.
Thus the above estimate shows that the averaged function $h_x$ is uniformly strictly concave.
\end{proof}

Let $h$ be the infimum of $h_x$ over all $x\in S(p,R)$, which is also $1$-Lipschitz and $(-\lambda)$-concave.
Furthermore, $h(p)\ge\sup_{x\in S(p,R)}h(x)+R/6$.
Set
\[U(t):=h^{-1}(a+t,\infty)\cap B(p,R),\quad a:=\sup_{x\in S(p,R)}h(x),\]
where $0\le t\ll R$.
Then $U(t)$ is convex.
Moreover, since $h$ is $1$-Lipschitz and $R\gg 1$, it follows that $U(t)$ contains $B(p,1)$.
Note that $U(t)$ also contains a small ball around the unique maximum point $q$ of $h$.

Finally we lift the situation to $\tilde M_i$.
Let $\tilde h_i$ and $\tilde U_i(t)$ denote the lifts of $h$ and $U(t)$ respectively.
Then $\tilde h_i$ is $1$-Lipschitz and $(-\lambda)$-concave on $\tilde B(p_i,2R)$, where the latter follows from the same argument using volume comparison as in Claim \ref{clm:conc} (note $\dim\tilde M_i=n$).
Therefore $\tilde U_i(t)$ is a convex domain containing $\tilde B(p_i,1)$.
Since the unique maximum point $q_i$ of $\tilde h_i$ converges to $q$, the condition (2) is clear.

From now on we fix $t$ and denote $U:=U(t)$ and $\tilde U_i:=\tilde U_i(t)$.

\step{Step 2}
Next we show that $U$ can be contracted to a tiny ball by the gradient flow of $h$ and this tiny ball can be crushed to the center by another gradient flow.
Then we check that $\tilde U_i$ can be contracted into a tiny ball in the same manner.
However, crushing this tiny ball needs a more careful argument, which is deferred to the final step.

As before, let $q\in U$ denote the unique maximum point of $h$.
For any $x\in U\setminus\{q\}$, by Lemma \ref{lem:grad} we have
\[dh(\uparrow_x^q)\ge\lambda|qx|/2.\]
In particular, for any fixed $r>0$, the function $h$ is $(\lambda r/4)$-regular in the direction to $q$ outside $B(q,r/2)$.
Hence Lemma \ref{lem:reg} implies that the gradient flow of $h$ pushes $U$ into $B(q,r/2)$ within time at most $8R/\lambda r$.
Since $h$ is $1$-Lipschitz and $0$-concave, this flow is $1$-Lipschitz (Proposition \ref{prop:lip}).
In conclusion,

\begin{clm}\label{clm:flow1}
The gradient flow of $h$ gives a $(1,8R/\lambda r)$-Lipschitz homotopy pushing $U$ into $B(q,r/2)$.
\end{clm}

We choose $r>0$ so small that ${(2r)}^{-1}S(q,2r)$ is sufficiently close to the space of directions $\Sigma_q$ at $q$ in the Gromov-Hausdorff distance.
In particular, there exists $\nu>0$ such that for any $x\in B(q,r)\setminus\{q\}$ there is $y\in S(q,2r)$ with $\tilde\angle qxy>\pi/2+\nu$ (actually $\nu\to\pi/2$ as $r\to 0$, but we cannot use such a strong regularity once the situation is lifted to $\tilde M_i$; see Claim \ref{clm:stab}).
This implies that the gradient flow of $\dist_{S(q,2r)}$ gives a Lipschitz strong contraction of $B(q,r)$ to $q$.
Indeed, for any $x\in B(q,r)\setminus\{q\}$, let $y\in S(q,2r)$ be as above and let $z\in S(q,2r)$ be closest to $x$.
Then we have
\[d\dist_{S(q,2r)}(\uparrow_x^q)\ge-\cos\tilde\angle zxq\ge-\cos\tilde\angle yxq\ge\sin\nu.\]
In other words, $\dist_{S(q,2r)}$ is $(\sin\nu)$-regular in the direction to $q$ on $B(q,r)$.
Hence Lemma \ref{lem:reg} implies that the gradient flow of $\dist_{S(q,2r)}$ crushes $B(q,r)$ to $q$ within time at most $r/\sin\nu$.
This flow is $c(r,\nu)$-Lipschitz on $B(q,r)$, where $c(r,\nu)$ is a constant depending only on $r$ and $\nu$ (Proposition \ref{prop:lip}; actually independent of $r$).
In conclusion,

\begin{clm}\label{clm:flow2}
The gradient flow of $\dist_{S(q,2r)}$ gives a $(c(r,\nu),r/\sin\nu)$-Lipschitz contraction of $B(q,r)$ to $q$.
\end{clm}

Combining the two gradient flows above, we get a $(C,1)$-Lipschitz strong contraction of $U$ to $q$ for some $C>0$.
We verify that the same construction works for $\tilde U_i$.
As before, let $q_i\in\tilde U_i$ denote the unique maximum point of $\tilde h_i$, which converges to $q$.
The same argument as Claim \ref{clm:flow1} using the $(-\lambda)$-concavity of $\tilde h_i$ shows that $\tilde U_i$ can be pushed into $\tilde B(q_i,r/2)$ by the gradient flow of $\tilde h_i$ with uniform Lipschitz constant and time.
However, in general $\tilde B(q_i,r)$ is not Lipschitz strongly contractible to $q_i$.
To contract this ball, we need to replace $q_i$ with another close point, which will be done in the final step.

\step{Step 3}
This part is the same as the proof of \cite[3.5]{MY:stab}, which is a noncollapsing version of \cite[3.2]{Y:ess}.
We show that there exists $\hat q_i\in\tilde M_i$ converging to $q$ such that the distance function from $\tilde S(\hat q_i,2r)$ is uniformly regular (independent of $i$) on $\tilde B(\hat q_i,r)$ in the direction to $\hat q_i$.
Then the same argument as Claim \ref{clm:flow2} shows that $\tilde B(\hat q_i,r)$ can be contracted to $\hat q_i$ with uniform Lipschitz constant and time.

The following construction is similar to that in Step 1, but here we use the space of directions $\Sigma_q$ instead of the ideal boundary $N(\infty)$.
Recall that $(2r)^{-1}S(q,2r)$ is close to $\Sigma_q$.
For fixed $0<\theta\ll 1$ we take a maximal $\theta r$-discrete set $\{x_\alpha\}_\alpha$ in $S(q,2r)$.
For each $\alpha$ we take a maximal $\delta r$-discrete set $\{x_{\alpha\beta}\}_{\beta=1}^{N_\beta}$ in $S(q,2r)\cap B(x_\alpha,2\theta r)$, where $\delta>0$ will be determined in Claim \ref{clm:stab}.
Then the Bishop-Gromov inequality implies
\begin{equation}\label{eq:low}
N_\alpha\ge c\left(\theta/\delta\right)^{n-1},
\end{equation}
where $c>0$ depends only on $n$, $D$, and $v$.

We define a function $f$ on $B(q,r)$ by
\[f:=\min_\alpha f_\alpha,\quad f_\alpha:=\frac1{N_\alpha}\sum_{\beta=1}^{N_\alpha}|x_{\alpha\beta}\cdot|.\]
Since $\{x_\alpha\}_\alpha$ is a $\theta r$-net, where $\theta\ll1$, it is easy to see that
\[f(q)\ge f(x)+|qx|/2\]
for any $x\in B(q,r)$.

Now we lift the situation to $\tilde M_i$.
Let $f_i$ denote the lift of $f$.
We define $\hat q_i$ as a maximum point of $f_i$ on $\tilde B(q_i,r)$, which converges to the unique maximum point $q$ of $f$.

\begin{clm}\label{clm:stab}
There exists $\nu>0$ such that for any sufficiently large $i$ and any $x_i\in\tilde B(\hat q_i,r)$, there is $y_i\in\tilde S(\hat q_i,2r)$ with $\tilde\angle\hat q_ix_iy_i>\pi/2+\nu$.
\end{clm}

\begin{proof}
Suppose the claim does not hold.
After passing to a subsequence, we can assume that there are $\nu_i\to 0$ and $x_i\in\tilde B(\hat q_i,r)$ such that $\tilde\angle\hat q_ix_iy_i<\pi/2+\nu_i$  for any $y_i\in\tilde S(\hat q_i,2r)$.
Note that $|\hat q_ix_i|\to 0$ since the corresponding claim holds on $N$ by the choice of $r$.
We may assume that $(|\hat q_ix_i|^{-1}\tilde M_i,\hat q_i)$ converges to a nonnegatively curved Alexandrov space $(\hat N,\hat q)$.

Let $x_{\alpha\beta}^i$ and $f_\alpha^i$ denote the lifts of $x_{\alpha\beta}$ and $f_\alpha^i$ defining $f_i$, respectively.
Let $v_i$ and $v_{\alpha\beta}^i$ denote the directions of shortest paths from $\hat q_i$ to $x_i$ and $x_{\alpha\beta}^i$, respectively (choose them so that the following first variation formula holds).
Fix $\alpha$ such that $df_i(v_i)=df_\alpha^i(v_i)$ for infinitely many $i$.
Since $f_i$ attains a maximum at $\hat q_i$, this implies
\[\frac1{N_\alpha}\sum_\beta-\cos\angle(v_{\alpha\beta}^i,v_i)\le 0.\]
We may assume that $v_i$ and $v_{\alpha\beta}^i$ converge to the directions at $\hat q$ of a shortest path and rays, respectively.
Then the lower semicontinuity of angle implies
\begin{equation}\label{eq:up}
\frac1{N_\alpha}\sum_\beta-\cos\angle(v_{\alpha\beta},v)\le 0.
\end{equation}

On the other hand, the assumption at the beginning of the proof implies that $\angle(v,v_{\alpha\beta})\ge\pi/2$.
Moreover, since $\{x_{\alpha\beta}\}_{\beta}$ is $\delta r$-discrete, we see that $\{v_{\alpha\beta}\}_\beta$ is $\delta/2$-discrete.
Thus a volume comparison argument using \eqref{eq:low} shows that for the majority of $\beta$, $\angle(v,v_{\alpha\beta})$ is uniformly bounded away from $\pi/2$ by a positive constant depending only on $c$, provided $\delta$ is small enough. 
However, this implies that the average of the cosines is strictly greater than $0$, which contradicts \eqref{eq:up}.
\end{proof}

As we have seen in Claim \ref{clm:flow2}, Claim \ref{clm:stab} is enough for the gradient flow of $\dist_{\tilde S(\hat q_i,2r)}$ to be a Lipschitz strong contraction of $\tilde B(\hat q_i,r)$ to $\hat q_i$ with uniform Lipschitz constant and time.
Thus its composition with the gradient flow of $\tilde h_i$ gives a Lipschitz weak contraction of $\tilde U_i$ to $\hat q_i$.
Furthermore, it is easy to modify the first flow to fix a small neighborhood of $q_i$.
For example, the following flow fixes $\tilde B(q_i,r/10)$ while keeping the Lipschitz constant and time uniformly bounded above:
\[\hat \Phi_i(x,t):=\Phi_i(x,\rho(|q_ix|)t),\]
where $\Phi_i$ is the gradient flow of $h_i$ and $\rho:[0,\infty)\to[0,1]$ is a $5/r$-Lipschitz function such that $\rho\equiv0$ on $[0,r/10]$ and $\rho\equiv1$ on $[r/5,\infty)$.
Since $\hat q_i$ also converges to $q$, this modified flow fixes $\hat q_i$ for sufficiently large $i$.
Therefore, its composition with the second flow gives a $(C,1)$-Lipschitz strong contraction of $\tilde U_i$ to $\hat q_i$ for some constant $C>0$ independent of $i$.
This completes the proof of Theorem \ref{thm:conv}.
\end{proof}

\begin{rem}
It is also possible to show that the distance to the center $\hat q_i$ is nonincreasing under the contraction (in other words, $\tilde U_i$ is \textit{strongly Lipschitz contractible} to $\hat q_i$ in the sense of \cite[1.4]{MY:stab}).
This is trivial for the second flow, but not for the first flow.
To see this for the first flow, we use the $(-\lambda)$-concavity and $1$-Lipschitz continuity of $\tilde h_i$ to obtain
\begin{align*}
d\tilde h_i(\uparrow_x^{\hat q_i})&\ge\frac{\tilde h_i(\hat q_i)-\tilde h_i(x)+\lambda|\hat q_ix|^2/2}{|\hat q_ix|}\\
&\ge\frac{\tilde h_i(q_i)-|q_i\hat q_i|-\tilde h_i(x)+\lambda|\hat q_ix|^2/2}{|\hat q_ix|}\ge-\frac{|q_i\hat q_i|}{|\hat q_ix|}+\lambda\frac{|\hat q_ix|}2.
\end{align*}
Since $|q_i\hat q_i|\to 0$ as $i\to\infty$, the last term is greater than $\lambda r/50$ outside $\tilde B(\hat q_i,r/20)$, which means that $h_i$ is $(\lambda r/50)$-regular in the direction to $\hat q_i$ outside $\tilde B(\hat q_i,r/20)$.
On the other hand, the modified flow $\tilde\Phi_i$ defined above fixes $\tilde B(\hat q_i,r/20)$.
Thus by Lemma \ref{lem:reg} the distance to $\hat q_i$ is nonincreasing under this modified flow.
\end{rem}

\section{Quantitative good coverings}\label{sec:good}

In this section we show the existence of a quantitative version of a good covering introduced in \cite{MY:good}.
Fix $n$, $D$, and $v$.

\begin{thm}\label{thm:good}
There is $C=C(n,D,v)>0$ satisfying the following:
for any $M\in \mathcal A(n,D,v)$ and $0<\epsilon<C^{-1}$, there exists a finite open covering $\mathcal U=\{ U_j\}_{j=1}^{N}$ of $M$ such that
\begin{enumerate}
\item any $\epsilon$-ball in $M$ is contained in some $U_j$;
\item for each $j$, the number of $U_i$ intersecting $U_j$ is at most $C$;
\item any intersection of $U_j$ is $(C,\epsilon)$-Lipschitz (strongly) contractible;
\item for any $M'\in\mathcal A(n,D,v)$ that is $(C^{-1}\epsilon)$-close to $M$, there is a lift $\mathcal U'$ of $\mathcal U$ to $M'$ that satisfies the same properties as (1)--(3) and has the same nerve as $\mathcal U$, i.e., a collection of $U_j$ has nonempty intersection if and only if so does its lift.
\end{enumerate}
\end{thm}

\begin{rem}
The covering number $N$ depends on $\epsilon$.
\end{rem}

Actually each element of $\mathcal U$ is a superlevel set of a strictly concave function given by Theorem \ref{thm:conv}.
Excluding the stability condition (4), we define

\begin{dfn}\label{dfn:good}
For a general metric space $M$, a (not necessarily finite) open covering $\mathcal U$ satisfying the conditions (1)--(3) is called a \textit{$(C,\epsilon)$-good covering}.
\end{dfn}

\begin{proof}[Proof of Theorem \ref{thm:good}]
To clarify the notation, we will denote by $C_0$ the constant of Theorem \ref{thm:conv}.
Take a maximal $\epsilon/2$-discrete set $\{p_j\}_{j=1}^N$ of $M$.
For each $j$ and $0\le t\le C_0^{-1}$, let
\[U_j(t)=\{h_j>-t\epsilon\}\]
be a $(C_0,\epsilon)$-Lipschitz contractible domain containing $B(p_j,2\epsilon)$, where $h_j$ is a strictly concave function as in Theorem \ref{thm:conv}(1).
Similarly we define its lift $U_j'(t)$ to $M'$ by using $h_j'$ as in Theorem \ref{thm:conv}(3).

We prove that there are $0\le t_j\le C_0^{-1}$ (depending on $j$) such that the coverings $\{U_j(t_j)\}_j$ and $\{U_j'(t_j)\}_j$ satisfy the conditions (1)--(4) above.
Actually the conditions (1) and (2) are satisfied for any choice of $t_j$.
Indeed, (1) follows from the properties of $p_j$ and $U_j(t)$ and (2) follows from these plus the Bishop-Gromov inequality.

To obtain (3) and (4), we show

\begin{clm}\label{clm:good}
For each $j$, there exists $0\le t_j\le C_0^{-1}$ such that
\begin{enumerate}
\item[(i)] any nonempty intersection of $U_j(t_j)$ contains the $(C_1^{-1}\epsilon)$-ball around its \textit{peak}, i.e., the unique maximum point of the strictly concave function defining the intersection, where $C_1>0$ depends only on $n$, $D$, and $v$;
\item[(ii)] the same holds for $U_j'(t_j)$, and moreover, $\{U_j'(t_j)\}_j$ has the same nerve as $\{U_j(t_j)\}_j$.
\end{enumerate}
\end{clm}

\begin{proof}
We prove it by induction on $j$.
For the base case, we put $t_1:=0$.
Indeed, $U_1(0)$ and $U_1'(0)$ contain the $(C_0^{-1}\epsilon)$-balls around their peaks, respectively, by Theorem \ref{thm:conv}(2).
To show the induction step, suppose there are $t_1,\dots,t_{j-1}$ such that the conclusion holds for $U_k(t_k)$ and $U_k'(t_k)$ ($1\le k\le j-1$).

\step{Step 1}
We define $t_j$ as follows.
Put $U_k:=U_k(t_k)$ for $1\le k\le j-1$ and consider a parametrized family $\mathcal U_j(t):=\{U_k,U_j(t)\}_{k=1}^{j-1}$ of subcoverings, where $0<t<C_0^{-1}$.
Let $\mathcal N_j(t)$ denote the nerve of $\mathcal U_j(t)$, which is nondecreasing in $t$ with respect to the inclusion relation.
Since the number of $U_k$ intersecting $U_j(C_0^{-1})$ is bounded above by a constant depending only on $n$, $D$, and $v$, one can find $0<t_j<C_0^{-1}$ such that
\begin{equation}\label{eq:N}
\mathcal N_j(t_j-c_0)=\mathcal N_j(t_j+c_0),
\end{equation}
where $c_0$ is a uniform positive constant independent of $\epsilon$.

\step{Step 2}
We first show (i).
Put $U_j:=U_j(t_j)$ and let $V$ be any intersection of $U_k$ ($1\le k\le j-1$) such that $V\cap U_j$ is nonempty.
Then the above property \eqref{eq:N} implies that $V\cap U_j^{-}$ is also nonempty, where $U_j^-:=U_j(t_j-c_0)$.
Let $x\in V\cap U_j^-$.

We observe that one can shrink $V$ and expand $U_j^-$ a little bit so that they still have nonempty intersection.
In view of Theorem \ref{thm:conv}(1), we may assume that $V=\{h>0\}$, where $h$ is a $1$-Lipschitz ($-(C_0\epsilon)^{-1}$)-concave function.
We set
\[V^-:=\{h>c_1\epsilon\},\quad\tilde U_j^-:=\{h_j>-(t_j-c_0/2)\epsilon\},\]
where $c_1$ is a sufficiently small uniform positive constant.

We show that $V^-\cap\tilde U_j^-$ is nonempty by moving $x$ if necessary by the gradient flow of $h$.
The induction hypothesis implies that $V$ contains the $(C_1^{-1}\epsilon)$-ball around its peak.
On the other hand, since $h$ is ($-(C_0\epsilon)^{-1}$)-concave, the gradient estimate (Lemma \ref{lem:grad}) shows that $|\nabla h|>(4C_0C_1)^{-1}$ outside the concentric ball of half radius.
Hence if $c_1$ is small enough, one can find $y\in V^-$ such that $|xy|<c_0\epsilon/2$.
Since $h_j$ is $1$-Lipschitz, this implies $y\in\tilde U_j^-$, as desired.

Now let $p$ be the peak of $V\cap U_j$.
Since $V^-\cap\tilde U_j^-$ is nonempty, it contains $p$.
Since $h$ and $h_j$ are $1$-Lipschitz, we see that $B(p,\tilde C_1^{-1}\epsilon)$ is contained in $V\cap U_j$ for sufficiently large $\tilde C_1\gg C_1$.

Finally, note that the constant $C_1$ of the statement gets worse in each induction step, but the number of the actual steps is at most the dimension of the nerve of the largest covering $\{U_j(C_0^{-1})\}_j$, which is uniformly bounded above.
This proves (i).

\step{Step 3}
Next we show (ii).
Suppose $M'$ is $(C^{-1}\epsilon)$-close to $M$, where $C$ is large enough.
Note that the lifted function $h'$ as in Theorem \ref{thm:conv}(3) is $(C^{-1}\epsilon)$-close to the original function $h$ via the approximation.
For a superlevel set of a strictly concave function on $M$, we will use the prime symbol $\ '$ to denote its lift to $M'$, like $U_j'$ for $U_j$.
Let $V$ be any nonempty intersection of $U_k$ ($1\le k\le j-1$).

The argument in Step 2 implies that if $U_j$ intersects $V$, then $U_j'\cap V'$ contains the $(C_1^{-1}\epsilon)$-ball around its peak.
Indeed, since $\tilde U_j^-$ intersects $V^-$, small expansions of $(\tilde U_j^-)'$ and $(V^-)'$ also have nonempty intersection.
By the $1$-Lipschitz continuity of the defining functions, we obtain the claim.

It remains to prove that if $U_j$ does not intersect $V$, then $U_j'\cap V'$ is also empty.
Suppose $U_j$ does not intersect $V$.
Then by the property \eqref{eq:N}, so does $U_j^+:=U_j(t_j+c_0)$.
Consider
\[V^+:=\{h>-c_2\epsilon\},\quad\tilde U_j^+:=\{h_j>-(t_j+c_0/2)\epsilon\},\]
where $c_2$ is a sufficiently small uniform positive constant (the former expansion makes sense in view of Theorem \ref{thm:conv}(1)).
If $c_2$ is small enough, then $\tilde U_j^+\cap V^+$ is also nonempty: otherwise, by using the induction hypothesis one can shrink $V^+$ and expand $\tilde U_j^+$ as in Step 2 to show that $U_j^+\cap V$ is nonempty, a contradiction.
Therefore $U_j'$ does not intersect $V'$, which completes the proof (note that the number of the induction steps is uniformly bounded above as in Step 2).
\end{proof}

To prove the condition (3), let $U$ be any nonempty intersection of $U_j(t_j)$ with peak $p$, which includes $B(p,C^{-1}\epsilon)$ for some constant $C>0$ by Claim \ref{clm:good}(1).
As observed above, the gradient estimate (Lemma \ref{lem:grad}) implies that $U$ admits a $(C,\epsilon)$-Lipschitz deformation retraction into $B(p,C^{-2}\epsilon)$.
Moreover, by Theorem \ref{thm:conv}, $B(p,C^{-2}\epsilon)$ is $(C,\epsilon)$-Lipschitz contractible in $B(p,C^{-1}\epsilon)$.
Hence $U$ is $(C,\epsilon)$-Lipschitz contractible.

The remaining properties in the condition (4) follows from Claim \ref{clm:good}(2) in the same way.
This completes the proof.
\end{proof}

\begin{rem}\label{rem:CAT}
Another important class of metric spaces having such good coverings is CAT spaces (more generally, CBA spaces with uniform injectivity radius) with uniform doubling constant (see Section \ref{sec:main2} for the doubling condition).
Indeed, if $X$ is a $C$-doubling CAT($1$) space, we take a maximal $\epsilon/2$-discrete set $\{p_j\}_j$ and set $U_j:=B(p_j,2\epsilon)$, where $0<\epsilon\ll C^{-1}$.
Then the same argument as above shows that this covering satisfies the conditions (1) and (2) of Theorem \ref{thm:good} (use the doubling condition instead of the Bishop-Gromov inequality).
Moreover, it satisfies the condition (3) since the CAT($1$) property implies that
\begin{itemize}
\item any metric ball of radius less than $\pi/2$ is convex;
\item shortest paths are unique in such a metric ball;
\item the contraction of the metric ball to any point along shortest paths is uniformly Lipschitz.
\end{itemize}
In particular, the argument in the next section can be applied to such CAT spaces.
\end{rem}

\section{First proof of Theorem \ref{thm:main}}\label{sec:main1}

In this section we give the first proof of Theorem \ref{thm:main} using good coverings.
Since we have the stability property as in Theorem \ref{thm:good}(4), it suffices to prove the following

\begin{thm}\label{thm:nerve}
For any $C>0$ there is $\tilde C>0$ satisfying the following:
Let $M$ be a metric space with a $(C,\epsilon)$-good covering $\mathcal U$ in the sense of Definition \ref{dfn:good}.
Then $M$ is $(\tilde C,\epsilon)$-Lipschitz homotopy equivalent to the $\epsilon$-geometric realization of the nerve of $\mathcal U$.
\end{thm}

Let us first explain our terminology.

The \textit{nerve} of a locally finite covering $\mathcal U=\{U_j\}_j$ is a simplicial complex with vertices $\{v_j\}_j$ such that a collection of vertices spans a simplex if and only if the corresponding $U_j$ have nonempty intersection.
We denote it by $\mathcal N=\mathcal N_\mathcal U$.
Note that the condition (2) of Theorem \ref{thm:good} implies that the nerve of a $(C,\epsilon)$-good covering has uniformly bounded local finiteness and in particular uniformly bounded dimension.

Let $K$ be a locally finite simplicial complex.
For simplicity we assume that $K$ has only finitely many vertices $\{v_j\}_{j=1}^N$, but this is not necessary because of the local finiteness.
For $\epsilon>0$, the \textit{$\epsilon$-geometric realization} $|K|=|K|_\epsilon$ of $K$ is defined as follows.
We identify $v_j$ with the $j$-th standard vector of norm $\epsilon$ in Euclidean space, i.e.,
\[v_j=(0,\dots,\overset j\epsilon,\dots,0)\in\mathbb R^N.\]
The set $|K|$ is a subset of the convex hull of $\{v_j\}_j$ corresponding to $K$.
More precisely, a function $\theta:\{v_j\}_j\to[0,1]$ defines an element $\sum_j\theta(v_j)v_j$ of $|K|$ if it satisfies
\begin{enumerate}
\item $\sum_j \theta(v_j)=1;$
\item $\supp\theta$ defines a simplex of $K$ (also denoted by $\supp\theta$).
\end{enumerate}
We consider the length metric on $|K|$ induced by the standard Euclidean metric on $\mathbb R^N$.
We identify a simplex $\sigma\in K$ with its closed realization in $|K|$ and denote it by the same symbol $\sigma$.

\begin{conv}\label{conv:C}
In the following proof, we will use the same symbol $C$ to denote various large positive constants.
Whenever this symbol appears, it means that there exists a constant $C$ such that the claim holds (possibly different from $C$ in Theorem \ref{thm:nerve}, but depending only on it as $\tilde C$).
\end{conv}

\begin{proof}[Proof of Theorem \ref{thm:nerve}]
The proof is along the same line as in \cite[\S3]{MY:lip}, where the second and third authors proved the qualitative version of this theorem (based on the idea of \cite[9.4.15]{Sha}).
In fact all we have to do is repeat the $\epsilon$-scaled version of the previous argument.
We will skip detailed calculations if they were already done in \cite{MY:lip}.

Let $\mathcal U=\{U_j\}_j$ be a $(C,\epsilon)$-good covering of $M$.
Let $\mathcal N=\mathcal N_{\mathcal U}$ denote its nerve and $|\mathcal N|=|\mathcal N|_\epsilon$ its $\epsilon$-geometric realization.
Note that the condition (2) of Theorem \ref{thm:good} implies that the dimension of $\mathcal N$ is at most $C$.
This fact will be used throughout the proof.
We prove that there exist $C$-Lipschitz maps
\[\Theta:M\to|\mathcal N|,\quad\zeta:|\mathcal N|\to M\]
such that $\zeta\circ\Theta$ and $\Theta\circ\zeta$ are $(C,\epsilon)$-Lipschitz homotopic to the identities of $M$ and $|\mathcal N|$, respectively.

\step{Step 1}
We first define the map $\Theta$ by using the partition of unity subordinate to $\mathcal U$.
Let $f_j$ be the distance function to the complement of $U_j$ and set
\[\xi_j(x):=\frac{f_j(x)}{\sum_if_i(x)}\]
for $x\in M$ (the above $f_j$ is simpler than the one in \cite{MY:lip}, but this is not an essential change).
Then

\begin{clm}\label{clm:xi}
$\xi_j$ is $C/\epsilon$-Lipschitz.
\end{clm}

\begin{proof}
By the condition (1) of Theorem \ref{thm:good}, we have $\sum_if_i(x)\ge\epsilon$ for any $x\in M$.
A direct calculation using the triangle inequality shows
\begin{align*}
&\left|\xi_j(x)-\xi_j(y)\right|\\
&\le\left|\frac{f_j(x)-f_j(y)}{\sum_if_i(x)}\right|+\left|\frac{\sum_if_i(x)-\sum_if_i(y)}{\sum_if_i(x)}\right|\frac{f_j(y)}{\sum_if_i(y)}.
\end{align*}
Since $f_j$ is $1$-Lipschitz, the first term is bounded above by $|xy|/\epsilon$.
Similarly, since $\dim\mathcal N\le C$, the second term is bounded above by $C|xy|/\epsilon$.
\end{proof}

We define $\Theta:M\to |\mathcal N|$ by
\[\Theta(x):=\sum_j\xi_j(x)v_j,\]
where $v_j$ is the $j$-th standard vector of norm $\epsilon$.
Then Claim \ref{clm:xi} together with $\dim\mathcal N\le C$ implies that $\Theta$ is $C$-Lipschitz.

\step{Step 2}
Next we construct larger spaces $\mathcal D$ and $\mathcal M$ with $C$-bi-Lipschitz embeddings
\[
\begin{CD}
M @>\tau >> \mathcal D\\
@.
@VV{\iota}V \\
|\mathcal N| @>>\Psi> \mathcal M
\end{CD}
\]
whose images are $(C,\epsilon)$-Lipschitz strong deformation retracts of the target spaces.
Only the deformation retraction for the $\iota$-image is nontrivial, whose construction is deferred to the next step.

For a simplex $\sigma\in\mathcal N$, we denote by $U_\sigma$ the corresponding intersection of $U_j$.
We define $\mathcal D$ by
\begin{align*}
\mathcal D&:=\{(\theta, x)\in |\mathcal N|\times M \,|\, x\in U_{\supp\theta}\} \\
 & =\bigcup_{\sigma\in\mathcal N} \sigma\times U_\sigma\subset |\mathcal N|\times M.
\end{align*}
As above we will use $\theta$ and $x$ to denote the $|\mathcal N|$- and $M$-coordinates, respectively.
Let $p:\mathcal D\to|\mathcal N|$ and $q:\mathcal D\to M$ be the projections.
The $C$-bi-Lipschitz embedding $\tau:M\to\mathcal D$ is given by
\[\tau(x):=(\Theta(x),x),\]
whose inverse is the restriction of $q$.

\begin{clm}\label{clm:tau}
$\tau(M)$ is a $(C,\epsilon)$-Lipschitz strong deformation retract of $\mathcal D$.
\end{clm}

\begin{proof}
The desired homotopy $H:\mathcal D\times[0,\epsilon]\to\mathcal D$ is given by
\[H(\theta,x,s):=((s/\epsilon)\Theta(x)+(1-s/\epsilon)\theta,x)\]
for $s\in[0,\epsilon]$.
Since each simplex has size $\epsilon$, this map is $C$-Lipschitz.
\end{proof}

We define $\mathcal M$ as the mapping cylinder of $p$ with height $\epsilon$:
\begin{align*}
\mathcal M=\mathcal M_\epsilon(p)&:= \mathcal D\times [0,\epsilon]\cup |\mathcal N|/(\theta, x,\epsilon)\sim \theta\\
&=\bigcup_{\sigma\in K} \sigma\times K(U_{{\sigma}})\subset |\mathcal N|\times K(M).
\end{align*}
Here $K(M)=K_\epsilon(M):=M\times [0,\epsilon]/M\times \epsilon$ denotes the Euclidean cone of height $\epsilon$ equipped with the metric
\begin{align*}
&\left|[x,t],[x',t']\right|^2\\
&:=(\epsilon- t)^2+(\epsilon-t')^2-2(\epsilon-t)(\epsilon-t')\cos\min\{\pi,|xx'|\},
\end{align*}
for $[x,t],[x',t']\in K(M)$ (the reason for this ``reversed'' metric is that the height corresponds to the time of homotopy later).
We consider the metric of $\mathcal M$ induced from the product metric of $|\mathcal N|\times K(M)$.

The maps $\iota:\mathcal D\to\mathcal M$ and $\Psi:|\mathcal N|\to\mathcal M$ are the natural isometric embeddings.
There is also a natural map $\Psi':\mathcal M\to|\mathcal N|$ induced by the projection $p$, whose restriction is the inverse of $\Psi$.
More precisely, 
\[\iota(\theta,x):=(\theta,x,0),\quad\Psi(\theta):=(\theta,v_M),\quad \Psi'(\theta,[x,t]):=\theta\]
where $v_M$ is the vertex of the cone $K(M)$.

Note that we have the following commutative diagram:

\begin{equation}\label{eq:Theta}
\begin{CD}
M @>\tau>> \mathcal D \\
@V\Theta VV @VV\iota V \\
|\mathcal N| @<<\Psi'< \mathcal M.
\end{CD}
\end{equation}

\begin{clm}\label{clm:Psi}
$\Psi(|\mathcal N|)$ is a $(C,\epsilon)$-Lipschitz strong deformation retract of $\mathcal M$.
\end{clm}

\begin{proof}
The desired homotopy $F:\mathcal M\times[0,\epsilon]\to\mathcal M$ is given by
\[F(\theta,[x,t],s):=(\theta, [x, (1-s/\epsilon)t+s])\]
for $s\in[0,\epsilon]$.
Since the cone has height $\epsilon$, this map is $C$-Lipschitz.
The details are left to the reader (see the proof of \cite[3.3]{MY:lip}).
\end{proof}

\step{Step 3}
Here we prove

\begin{clm}\label{clm:Phi}
There exists a $(C,\epsilon)$-Lipschitz strong deformation retraction $\Phi$ of $\mathcal M$ to $\iota(\mathcal D)$.
\end{clm}

We will construct $\Phi$ by reverse induction on skeleta.
As in \cite[\S3]{MY:lip}, since the geometric realization has the length metric, the problem is reduced to the following claim on each simplex.

\begin{subclm}\label{subclm:Phi}
For any $\sigma\in\mathcal N$, there exists a $(C,\epsilon)$-Lipschitz strong deformation retraction $\Phi_\sigma$ of $\sigma\times K(U_{\sigma})$ to $\sigma\times U_{\sigma}\times 0\cup\partial\sigma\times K(U_{\sigma})$.
\end{subclm}

Indeed, the desired map $\Phi$ is obtained by gluing $\Phi_\sigma$ by reverse induction on skeleta.
Since $\dim\mathcal N\le C$, the Lipschitz constant and the time of homotopy are uniformly bounded above.

\begin{rem}
There was a minor mistake in the proof of the corresponding claim in \cite{MY:lip} which might confuse the reader.
In the proofs of Sublemma 3.6 and Claim 3.8 in \cite{MY:lip}, we assumed that the the Lipschitz contraction $\varphi$ satisfies $\varphi(\cdot,t)\equiv p$ for $t\ge L/2$.
However, this choice of a constant was not sufficient for the proofs.
A suitable choice is, for example, $\varphi\equiv p$ for $t\ge L/10$.
Actually one may choose an arbitrary small constant $c>0$ so that $\varphi\equiv p$ for $t\ge cL$.
\end{rem}

\begin{proof}
The basic idea is that we first contract a neighborhood of the vertex of the cone, then the rest is almost a product, on which Lipschitz continuity is easy to prove.

To simplify the notation we will omit the subscript $\sigma$.
Let $\varphi:U\times[0,\epsilon] \to U$ be a $(C,\epsilon)$-Lipschitz contraction with center $p\in U$.
We may assume that
\begin{equation}\label{eq:varphi}
\varphi(\cdot,t)\equiv p
\end{equation}
for all $t\ge\epsilon/10$.
We divide the problem into two cases, depending on the dimension of $\sigma$.

\step{Case 1}
Suppose $\dim\sigma=0$.
Define a retraction $r:K(U)\to U\times 0$ by
\[r([x,t]):=[\varphi(x,t),0],\]
which is well-defined for $t=\epsilon$.

Let $g:[0,\epsilon]\to[0,\epsilon]$ be a $C$-Lipschitz function such that $g\equiv\epsilon$ on $[0,\epsilon/3]$ and $g(\epsilon)=0$.
Define $\Phi:K(U)\times[0,\epsilon]\to K(U)$ by
\[\Phi([x,t],s):=[\varphi(x,(s/\epsilon)t),(g(s)/\epsilon)t]\]
for $s\in[0,\epsilon]$.
Then it gives a $(C,\epsilon)$-Lipschitz homotopy between the identity and $r$.
The details are left to the reader (see the proof of \cite[3.8]{MY:lip}).

\step{Case 2}
Suppose $\dim\sigma\ge 1$.
Let
\[\rho=(\psi,u):\sigma\times [0,\epsilon]\to\sigma\times 0\cup\partial\sigma\times [0,\epsilon]\]
be the radial projection from $(\theta_*,2\epsilon)$, where $\theta_*$ is the barycenter of $\sigma$. 
Since $\sigma$ is a simplex of size $\epsilon$, $\rho$ is a $C$-Lipschitz retraction.
Define a retraction $r:\sigma\times K(U)\to\sigma\times U\times0\cup\partial\sigma\times K(U)$ by
\begin{equation}\label{eq:r}
r(\theta,[x,t]):=(\psi(\theta,t),[\varphi(x,t-u(\theta,t)),w(\theta,t)]),
\end{equation}
where $w:\sigma\times[0,\epsilon]\to[0,\epsilon]$ is defined as follows.
Set
\[\Omega_0:=u^{-1}[0,\epsilon/10],\quad\Omega_1:=u^{-1}[\epsilon/2,\epsilon]\]
and let $s_i$ be the distance functions to $\Omega_i$ ($i=0,1$).
We define $w$ by 
\[w(\theta,t):=\frac{s_1}{s_0+s_1}u+\frac{s_0}{s_0+s_1}t.\]
Then $w=u$ on $\Omega_0$ and $w=t$ on $\Omega_1$.
Since $\sigma$ is of size $\epsilon$, $w$ is $C$-Lipschitz.
Note that $r$ is well-defined for $t=\epsilon$.

Let $\mu$ and $\nu$ be $C$-Lipschitz functions from $[0,\epsilon]$ to $[0,\epsilon]$ such that
\begin{equation*}
\mu(s)=
\begin{cases}
0 & s\le\epsilon/2\\
\epsilon & s\ge2\epsilon/3,
\end{cases}
\quad
\nu(s)=
\begin{cases}
0 & s\le2\epsilon/3\\
\epsilon & s\ge3\epsilon/4.
\end{cases}
\end{equation*}
Define $\Phi:\sigma\times K(U)\times[0,\epsilon]\to\sigma\times K(U)$ by
\begin{align*}
\Phi(\theta,[x,t],s):=((&1-(s/\epsilon))\theta+(s/\epsilon)\psi,\\
&[\varphi(x,(\mu/\epsilon)(t-u)),((1-(\nu/\epsilon))t+(\nu/\epsilon)w)])
\end{align*}
for $s\in[0,\epsilon]$.
Then it gives a $(C,\epsilon)$-Lipschitz homotopy between the identity and $r$.
The details are left to the reader (see the proof of \cite[3.6]{MY:lip}).
\end{proof}

\step{Step 4}
The inverse map $\zeta:|\mathcal N|\to M$ is given by the following commutative diagram:
\begin{equation}\label{eq:zeta}
\begin{CD}
M @<q << \mathcal D \\
@A\zeta AA @AA\bar\Phi A \\
|\mathcal N| @>>\Psi> \mathcal M,
\end{CD}
\end{equation}
where $\bar\Phi:=\Phi(\cdot,\epsilon)$ and $\mathcal D$ is identified with its $\iota$-image.
Now the statement follows from \eqref{eq:Theta}, \eqref{eq:zeta}, Claims \ref{clm:tau}, \ref{clm:Psi}, and \ref{clm:Phi}.
\end{proof}

\begin{rem}\label{rem:nerve}
By construction, for any $x\in M$ we have $x\in U_{\supp\Theta(x)}$.
On the other hand, for any $\theta\in|\mathcal N|$ there exists a vertex of $\supp\theta$ such that the corresponding $U_j$ contains $\zeta(\theta)$.
Therefore $\zeta\circ\Theta(x)$ is contained in some $U_j$ containing $x$.
\end{rem}

\begin{proof}[Proof of Theorem \ref{thm:main}]
Suppose $M,M'\in\mathcal A(n,D,v)$ are $(C^{-1}\epsilon)$-close in the Gromov-Hausdorff distance, where $C>0$ is a sufficiently large constant and $0<\epsilon<C^{-1}$.
By Theorem \ref{thm:good}, there exist a $(C,\epsilon)$-good covering $\mathcal U$ of $M$ and a corresponding $(C,\epsilon)$-good covering $\mathcal U'$ of $M'$ with the same nerve as $\mathcal U$.
By Theorem \ref{thm:nerve}, $M$ and $M'$ are $(C,\epsilon)$-Lipschitz homotopy equivalent to the $\epsilon$-geometric realization of the nerve.
Therefore $M$ and $M'$ are $(C,\epsilon)$-Lipschitz homotopy equivalent.
Furthermore, since $\mathcal U'$ is a lift of $\mathcal U$, an observation similar to Remark \ref{rem:nerve} shows that the obtained Lipschitz homotopy equivalences are $C\epsilon$-close to the original Gromov-Hausdorff approximations.
In particular they are also $C\epsilon$-approximations.
This completes the proof.
\end{proof}

\section{Second proof of Theorem \ref{thm:main}}\label{sec:main2}

In this section we give the second proof of Theorem \ref{thm:main} using only local contractibility.
We prove the following Lipschitz version of Petersen's stability theorem \cite[Theorem A]{Peter}.

\begin{thm}\label{thm:Peter}
For any $C>0$ there is $\tilde C>0$ satisfying the following:
Let $M$ and $M'$ be metric spaces that are locally $C$-Lipschitz contractible and $C$-doubling.
If their Gromov-Hausdorff distance is less than $\tilde C^{-1}$, then we have
\[d_{\tilde C\mathchar`-\mathrm{LH}}(M,M')<\tilde Cd_\mathrm{GH}(M,M').\]
More precisely, for any given $\epsilon$-approximation between $M$ and $M'$, where $\epsilon<\tilde C^{-1}$, there exists a $(\tilde C,\tilde C\epsilon)$-Lipschitz homotopy approximation between $M$ and $M'$ that is $\tilde C\epsilon$-close to it.
\end{thm}

Let us first explain our terminology.
Fix $C>0$.
A metric space is \textit{locally $C$-Lipschitz contractible} if any $\epsilon$-ball is contained in some $(C,\epsilon)$-Lipschitz contractible domain, where $0<\epsilon<C^{-1}$.
A metric space is \textit{$C$-doubling} if any $\epsilon$-ball is covered by at most $C$ balls of radius $\epsilon/2$, where $0<\epsilon<C^{-1}$.

For Alexandrov spaces, the former condition is guaranteed by Theorem \ref{thm:conv} and the latter one is a consequence of the Bishop-Gromov inequality.
Therefore Theorem \ref{thm:main} immediately follows from Theorem \ref{thm:Peter}.
As in Remark \ref{rem:CAT}, CAT spaces with a uniform doubling constant also satisfy the assumption of Theorem \ref{thm:Peter}.

The proof of Theorem \ref{thm:Peter} is along almost the same line as that of \cite[Theorem A]{Peter}.
Only the existence of a domination by a polyhedron (Proposition \ref{prop:dom}), which was used in \cite{Peter} without proof, will be proved by an argument as in the previous section.

In the following proof, we will use Convention \ref{conv:C} as in the previous section.
Furthermore we will use

\begin{conv}\label{conv:C'}
In this section $\tilde C$ denotes a large positive constant depending only on $C$ and $\epsilon$ denotes a small positive number less than some constant depending only on $C$.
More precisely, whenever these symbols appear, it means:
\begin{quote}
for any $C>0$ there exists $\tilde C>0$ such that for any $0<\epsilon<\tilde C^{-1}$ the following holds.
\end{quote}
To be precise, all the statements in this section should be preceded by this sentence, but we will omit it for simplicity.
\end{conv}

\subsection{Lipschitz maps from polyhedra}

We will use the same notation and terminology as in the previous section.
For a (locally finite) simplicial complex $K$, we denote by $|K|=|K_\epsilon|$ its $\epsilon$-geometric realization.
We identify a simplex $\sigma\in K$ with its closed realization in $|K|$ and denote it by $\sigma$.

Furthermore, if $L$ is a subcomplex $K$, we also identify $|L|$ with a subset of $|K|$ (equipped with the metric of $|K|$).
We denote by $K^0$ the $0$-skeleton of $K$, i.e., the set of vertices of $K$.

We first prove the Lipschitz version of \cite[\S2, Main Lemma]{Peter}.
The proof is based on an idea similar to that of Subclaim \ref{subclm:Phi}.

\begin{prop}\label{prop:simp}
Let $M$ be a locally $C$-Lipschitz contractible space and $K$ a simplicial complex of dimension $\le C$ with subcomplex $L$.
Let $|K|=|K|_\epsilon$ denote the $\epsilon$-geometric realization.
Suppose $f:|K^0|\cup|L|\to M$ is $C$-Lipschitz.
If $\epsilon$ is small enough, then there exists a $\tilde C$-Lipschitz extension $\tilde f$ of $f$ onto $|K|$.
\end{prop}

\begin{proof}
The proof is almost the same as that of \cite[\S2, Main Lemma]{Peter}.
However, we need to be more careful about Lipschitz constants.
Since the geometric realization has the length metric, it suffices to check the Lipschitz continuity on each simplex.

The map $\tilde f$ is constructed by induction on skeleta.
On $|K^0|\cup|L|$ we define $\tilde f:=f$, which includes the base case.
Let $\sigma\in K\setminus(K^0\cup L)$ be a simplex.
By the induction hypothesis we may assume that $f$ is already defined on $\partial\sigma$.
By the Lipschitz continuity the diameter of $f(\partial\sigma)$ is bounded above by $C\epsilon$.
Hence if $\epsilon$ is small enough, the assumption implies that there exists a $(C,\epsilon)$-Lipschitz contractible domain $U$ containing $f(\partial\sigma)$.
We may assume that the $(C,\epsilon)$-Lipschitz contraction $\varphi:U\times[0,\epsilon]\to U$ satisfies $\varphi(\cdot,t)\equiv p$ for all $t\ge\epsilon/2$, where $p$ is the center of $U$.

Let $x_\ast$ be the barycenter of $\sigma$.
Any point $x\in\sigma\setminus\{x_\ast\}$ is uniquely represented as a convex combination $t(x)x_\ast+(1-t(x))y(x)$, where $y(x)\in\partial\sigma$ and $0\le t(x)<1$ (we define $t(x_\ast):=1$).
Observe that the map $x\mapsto(y(x),\epsilon t(x))$ is a $C$-bi-Lipschitz homeomorphism between $t^{-1}[0,1/2]$ and the product $\partial\sigma\times[0,\epsilon/2]$.
Therefore if we define $\tilde f$ on $\sigma$ by
\[\tilde f(x):=\varphi(f\circ y(x),\epsilon t(x)),\]
then it gives a $C$-Lipschitz extension of $f$ onto $\sigma$ (well-defined for $x_\ast$).
Note that the Lipschitz constant will get worse at each induction step, but the number of the induction steps is uniformly bounded above by assumption.
\end{proof}

\begin{cor}\label{cor:simp}
Let $M$ be a locally $C$-Lipschitz contractible space and $K$ a simplicial complex of dimension $\le C$.
Let $|K|=|K|_\epsilon$ denote the $\epsilon$-geometric realization.
Suppose $C$-Lipschitz maps $f_i:|K|\to M$ ($i=0,\epsilon$) are $\epsilon$-close in the uniform distance.
If $\epsilon$ is small enough, then there exists a $(\tilde C,\epsilon)$-Lipschitz homotopy between them.
\end{cor}

\begin{proof}
Observe that $|K|\times[0,\epsilon]$ has a natural triangulation whose vertices are exactly those of $|K|\times\{0,\epsilon\}$ and such that $|K|\times[0,\epsilon]$ is $C$-bi-Lipschitz homeomorphic to the $\epsilon$-geometric realization of this triangulation.
The statement follows from Proposition \ref{prop:simp} by regarding $|K|\times\{0,\epsilon\}$ as $|L|$.
Note that the assumption that $f_0$ and $f_\epsilon$ are $\epsilon$-close implies that the glued map $f_0\cup f_\epsilon$ defined on $|K|\times\{0,\epsilon\}$ is $C$-Lipschitz.
\end{proof}

\subsection{Lipschitz dominations by polyhedra}\label{sec:dom}

The contents of this subsection have no counterparts in \cite{Peter}.

Let $M$ be a metric space that is locally $C$-Lipschitz contractible and $C$-doubling.
For sufficiently small $\epsilon>0$, take a maximal $\epsilon/2$-discrete set $\{p_j\}_j$ in $M$.
Instead of a good covering, we will consider an open covering $\mathcal U=\{B(p_j,2\epsilon)\}_j$.
Let $\mathcal N=\mathcal N_\mathcal U$ denote its nerve and $|\mathcal N|=|\mathcal N|_\epsilon$ its $\epsilon$-geometric realization.

By the choice of $p_j$, any $\epsilon$-ball in $M$ is contained in some $B(p_j,2\epsilon)$.
Furthermore, the $C$-doubling condition implies that $\mathcal N$ has uniformly bounded local finiteness.
In other words, this covering still satisfies the conditions (1) and (2) of Theorem \ref{thm:good} (but not (3)).

In this subsection we prove

\begin{prop}\label{prop:dom}
$M$ is $(\tilde C,\epsilon)$-Lipschitz dominated by $|\mathcal N|$, i.e., there exist $\tilde C$-Lipschitz maps $\Theta:M\to |\mathcal N|$ and $\zeta:|\mathcal N|\to M$ such that $\zeta\circ\Theta$ is $(\tilde C,\epsilon)$-Lipschitz homotopic to $1_M$.
\end{prop}

\begin{rem}
Unlike the case of a good covering, $\Theta\circ\zeta$ is not homotopic to $1_{|\mathcal N|}$ because of the lack of the condition (3) of Theorem \ref{thm:good} (see also Remark \ref{rem:tau'}).
\end{rem}

\begin{proof}
The proof is along almost the same line as that of Theorem \ref{thm:nerve}.
Instead of considering intersections of metric balls defined by $\mathcal N$, we will construct a family of $(C,\epsilon)$-Lipschitz contractible domains corresponding to $\mathcal N$, which enables us to prove only the one-sided homotopy equivalence.
Once such a family is constructed, the rest of the proof is the same as the previous one, so we will omit it.

\step{Step 1}
As before, the map $\Theta$ is defined by using the partition of unity subordinate to $\mathcal U$.
Let $f_j$ be the distance function to the complement of $B(p_j,\epsilon)$ and set
\[\xi_j(x):=\frac{f_j(x)}{\sum_i f_i(x)}\]
for $x\in M$.
Since $\mathcal U$ still satisfies the conditions (1) and (2) of Theorem \ref{thm:good}, the same argument as in the proof of Claim \ref{clm:xi} shows

\begin{clm}\label{clm:xi'}
$\xi_j$ is $C/\epsilon$-Lipschitz.
\end{clm}

We define the $C$-Lipschitz map $\Theta:M\to |\mathcal N|$ by $\Theta(x):=\sum_j\xi_j(x)v_j$, where $v_j$ is the $j$-th standard vector of norm $\epsilon$.

\step{Step 2}
Next we construct larger spaces $\mathcal D$ and $\mathcal M$ with $C$-bi-Lipschitz embeddings
\[
\begin{CD}
M @>\tau >> \mathcal D\\
@.
@VV{\iota}V \\
|\mathcal N| @>>\Psi> \mathcal M.
\end{CD}
\]
This is the only place where an essential modification is required.

We first assign to each simplex $\sigma\in\mathcal N$ a $(C,\epsilon)$-Lipschitz contractible domain $U_\sigma$ in $M$.
This is done by reverse induction on skeleta.
Let $\sigma\in\mathcal N$ be a maximal simplex (with respect to the inclusion relation).
We denote by $B_\sigma$ the corresponding nonempty intersection of $2\epsilon$-balls in $M$.
By the local $C$-Lipschitz contractibility of $M$, there exists a $(C,\epsilon)$-Lipschitz contractible domain $U_\sigma$ containing $B_\sigma$.
Next let $\tau\in\mathcal N$ be a nonmaximal simplex.
Then there exists a $(C,\epsilon)$-Lipschitz contractible domain $U_\tau$ containing $B_\tau$ and all $U_\sigma$ such that $\sigma$ has $\tau$ as a face, where $U_\sigma$ was obtained by the induction hypothesis.
The Lipschitz constant will get worse at each induction step, but the number of the induction steps is uniformly bounded above by the $C$-doubling condition.
In this way we obtain a family $\{U_\sigma\}_{\sigma\in\mathcal N}$ of $(C,\epsilon)$-Lipschitz contractible domains such that
\begin{enumerate}
\item $U_\sigma$ contains $B_\sigma$;
\item if $\tau$ is a face of $\sigma$, then $U_\tau$ contains $U_\sigma$.
\end{enumerate}

By using this family, we define $\mathcal D$ in the same way as before:
\begin{align*}
\mathcal D&:=\{(\theta, x)\in |\mathcal N|\times M \,|\, x\in U_{\supp\theta}\} \\
 & =\bigcup_{\sigma\in\mathcal N} \sigma\times U_\sigma\subset |\mathcal N|\times M.
\end{align*}
Here the second equality follows from the condition (2) above.
Let $p:\mathcal D\to|\mathcal N|$ and $q:\mathcal D\to M$ be the projections.
The $C$-bi-Lipschitz embedding $\tau:M\to\mathcal D$ is given by $\tau(x):=(\Theta(x),x)$ with inverse $q$.
Note that the image of $\tau$ is certainly contained in $\mathcal D$ because of the condition (1) above.
As before, we define $\mathcal M$ as the mapping cylinder of $p$ with height $\epsilon$.

The maps $\iota:\mathcal D\to\mathcal M$ and $\Psi:|\mathcal N|\to\mathcal M$ are natural isometric embeddings and $\Psi':\mathcal M\to|\mathcal N|$ is a map induced by the projection $p$.
Then we have the following commutative diagram:

\begin{equation}\label{eq:Theta'}
\begin{CD}
M @>\tau>> \mathcal D \\
@V\Theta VV @VV\iota V \\
|\mathcal N| @<<\Psi'< \mathcal M.
\end{CD}
\end{equation}

The following claim is proved in exactly the same way as Claim \ref{clm:Psi}.

\begin{clm}\label{clm:Psi'}
$\Psi(|\mathcal N|)$ is a $(C,\epsilon)$-Lipschitz strong deformation retract of $\mathcal M$.
\end{clm}

\begin{rem}\label{rem:tau'}
We cannot prove Claim \ref{clm:tau}: $\tau(M)$ is not necessarily a deformation retract of $\mathcal D$.
This is because $U_\sigma$ is no longer the intersection of $U_j$.
For this reason we have only the one-sided homotopy equivalence in Proposition \ref{prop:dom}.
\end{rem}

\step{Step 3}
The next claim is also proved in the same way as Claim \ref{clm:Phi}.

\begin{clm}\label{clm:Phi'}
There exists a $C$-Lipschitz retraction $\bar\Phi$ of $\mathcal M$ to $\iota(\mathcal D)$.
\end{clm}

\begin{rem}
In fact there exists a $(C,\epsilon)$-Lipschitz strong deformation retraction $\Phi$ such that $\bar\Phi=\Phi(\cdot,\epsilon)$ as in Claim \ref{clm:Phi}.
However, we do not need this homotopy here.
\end{rem}

The proof is by reverse induction on skeleta.
The desired map $\bar\Phi$ is obtained as the compositions of $\bar\Phi_\sigma$ defined on each simplex $\sigma\in\mathcal N$.
Note that the condition (2) in Step 2 guarantees that the composition can be defined.
The construction of $\bar\Phi_\sigma=\Phi_\sigma(\cdot,\epsilon)$ is exactly the same as in Subclaim \ref{subclm:Phi}, so we omit the proof.

\step{Step 4}
The inverse map $\zeta:|\mathcal N|\to M$ is defined by the following commutative diagram:
\begin{equation}\label{eq:zeta'}
\begin{CD}
M @<q << \mathcal D \\
@A\zeta AA @AA\bar\Phi A \\
|\mathcal N| @>>\Psi> \mathcal M.
\end{CD}
\end{equation}
Now the statement follows from \eqref{eq:Theta'}, \eqref{eq:zeta'}, Claims \ref{clm:Psi'}, and \ref{clm:Phi'}.
\end{proof}

\subsection{Lipschitz homotopy}

By using dominations, we obtain

\begin{prop}\label{prop:hom}
Let $M$ and $M'$ be metric spaces that are locally $C$-Lipschitz contractible and $C$-doubling.
Suppose $C$-Lipschitz maps $f_i:M\to M'$ ($i=0,\epsilon$) are $\epsilon$-close in the uniform distance.
If $\epsilon$ is small enough, then there exists a $(\tilde C,\epsilon)$-Lipschitz homotopy between them.
\end{prop}

\begin{proof}
Let $\Theta:M\to|\mathcal N|$ and $\zeta:|\mathcal N|\to M$ be $(C,\epsilon)$-Lipschitz dominations as in Proposition \ref{prop:dom}.
Then $f_i$ is $(C,\epsilon)$-Lipschitz homotopic to $f_i\circ\zeta\circ\Theta$.
Furthermore, Corollary \ref{cor:simp} together with the assumption implies that $f_i\circ\zeta$ ($i=0,\epsilon$) are $(C,\epsilon)$-Lipschitz homotopic.
\end{proof}

Now we are ready to prove Theorem \ref{thm:Peter}

\begin{proof}[Proof of Theorem \ref{thm:Peter}]
Let $M$ and $M'$ be as in the assumption that are $\epsilon$-close in the Gromov-Hausdorff distance, where $\epsilon$ is small enough.
Take a maximal $\epsilon/2$-discrete set $\{p_j\}_j$ of $M$ and construct a $C$-Lipschitz map $\Theta:M\to |\mathcal N|$ as in the proof of Proposition \ref{prop:dom}.
Choose a corresponding (not necessarily $\epsilon/2$-discrete) set $\{p_j'\}_j$ of $M'$ by the $\epsilon$-approximation and construct a $C$-Lipschitz map $f:|\mathcal N|\to M'$ by Proposition \ref{prop:simp} such that $f(v_j)=p_j'$ for all $j$, where $v_j$ is the $j$-th vertex corresponding to $B(p_j,\epsilon)$ (observe that $f$ is $C$-Lipschitz on $|\mathcal N^0|$).
Then $F:=f\circ\Theta$ is a $C$-Lipschitz map from $M$ to $M'$.
Furthermore, since $F$ maps $p_j$ to a point $C\epsilon$-close to $p_j'$, it is $C\epsilon$-close to the original $\epsilon$-approximation.
Similarly one can construct a $C$-Lipschitz map $F'$ from $M'$ to $M$ with the same property.
Since $F$ and $F'$ are based on the same approximation, $F'\circ F$ and $F\circ F'$ are $C\epsilon$-close to $1_M$ and $1_{M'}$, respectively.
Then Proposition \ref{prop:hom} shows that they are $(C,C\epsilon)$-Lipschitz homotopic, which completes the proof.
\end{proof}

\section{Proof of Theorem \ref{thm:isom}}\label{sec:isom}

In this section we prove Theorem \ref{thm:isom}.
Although the proof is essentially the same as in \cite[1.4]{MY:lip}, to clarify the Lipschitz constant estimate, we shall take a slightly different approach.

As in the introduction, we will use the notation $\mathcal R_M(\delta,\ell)$ and $\tau(\delta)$, where $\delta$ is less than a constant depending only on $n$.
In the following proofs we will abuse $\tau(\delta)$ to denote different functions.
We will also abuse $C$ to denote various positive constants, which depends only on $n$, $D$, and $v$ unless otherwise stated.

The following theorem was proved in \cite[9.8]{BGP} (cf.\ \cite{WSS}, \cite{Y:conv}) by using the center of mass technique.

\begin{thm}\label{thm:reg}
There exists $C=C(n)>0$ satisfying the following.
Let $M,M'$ be $n$-dimensional Alexandrov spaces with Gromov-Hausdorff distance $\epsilon<\delta\ell$.
Then there exists a $\tau(\delta)$-almost isometric open embedding of $\mathcal R_M(\delta,\ell)$ into $M'$, which is $C\epsilon$-close to the original $\epsilon$-approximation.
\end{thm}

We will glue the above almost isometry with the Lipschitz homotopy approximation constructed in Theorem \ref{thm:main}.
The next proposition is essential in our gluing argument, which was also used in \cite{MY:lip}.
Here we will omit the subscript $M$ of $\mathcal R_M(\delta,\ell)$ to simplify the notation.

\begin{prop}\label{prop:diag}
There exists $C>0$ (independent of $n$, $\delta$, and $\ell$) satisfying the following.
Let $M$ be an $n$-dimensional Alexandrov space.
Let $\Delta$ denote the diagonal set in $M\times M$.
Then for any $0<\epsilon<\delta\ell$, $\mathcal R(\delta,\ell)^2\cap B(\Delta,\epsilon)$ admits a $(C,\epsilon)$-Lipschitz deformation retraction into $\Delta$ in $B(\Delta,\epsilon)$, i.e., there exists a $C$-Lipschitz map\[\mathbf\Psi:\mathcal R(\delta,\ell)^2\cap B(\Delta,\epsilon)\times[0,\epsilon]\to B(\Delta,\epsilon)\]
such that $\mathbf\Psi(\cdot,\epsilon)\in\Delta$ and $\mathbf\Psi(\cdot,t)|_\Delta=1_\Delta$ for any $0\le t\le\epsilon$.
\end{prop}

The following proof is slightly different from the previous one in \cite{MY:lip} to make the Lipschitz constant estimate clearer.

\begin{proof}
The desired deformation retraction is given by modifying the gradient flow of the distance function from $S(\Delta,2\epsilon)$, which we will denote by $f$.
We first observe that $f$ is $(1-\tau(\delta))$-regular on $\mathcal R(\delta,\ell)^2\cap B(\Delta,\epsilon)\setminus\Delta$ in the direction to $\Delta$, i.e.,
\[df(\uparrow_{\mathbf x}^\Delta)>1-\tau(\delta)\]
for any $\mathbf x\in\mathcal R(\delta,\ell)^2\cap B(\Delta,\epsilon)\setminus\Delta$.

Let us write $\mathbf x=(x_1,x_2)$, where $x_i\in\mathcal R(\delta,\ell)$ and $0<|x_1x_2|<\sqrt 2\epsilon$.
Note that a point of $\Delta$ closest to $\mathbf x$ can be written as $\mathbf y=(y,y)$, where $y$ is a midpoint of a shortest path $x_1x_2$.
Assume you can extend $x_1x_2$ beyond both endpoints by the same length to a shortest path $z_1z_2$ of length $2\sqrt 2\epsilon$.
Then it is easy to see that $\mathbf z:=(z_1,z_2)$ is the unique closest point to $\mathbf x$ in $S(\Delta,2\epsilon)$ (the uniqueness follows from the fact that there are no branching geodesics in Alexandrov spaces).
In this case $f$ is $1$-regular at $\mathbf x$ in the direction to $\Delta$.
In general, by using a strainer, one can find an ``almost geodesic extension'' $z_1x_1x_2z_2$ in the sense that $\tilde\angle yx_iz_i>\pi-\tau(\delta)$, which is almost unique.
Thus we obtain the $(1-\tau(\delta))$-regularity of $f$.
The details are left to the reader.

Let $\mathbf\Phi$ be the gradient flow of $f$.
For $\mathbf x\in\mathcal R(\delta,\ell)^2\cap B(\Delta,\epsilon)$, we denote by $T(\mathbf x)$ the minimum time for which $\mathbf\Phi(\mathbf x,t)$ reaches $\Delta$.
The above observation, together with Lemma \ref{lem:reg}, shows $T(\mathbf x)<2\epsilon$.
Note that $\mathbf\Phi$ is uniformly $C$-Lipschitz (independent of any a priori constant) for time $<2\epsilon$.
This is a consequence of Proposition \ref{prop:lip} and the fact that $f$ is $(2\epsilon^{-1})$-concave on $B(\Delta,\epsilon)$ (the latter follows from triangle comparison; see \cite[2.4]{MY:loc}).

Furthermore, the map $\mathbf x\mapsto T(\mathbf x)$ is $2C$-Lipschitz.
Indeed, let $\mathbf y\in\mathcal R(\delta,\ell)^2\cap B(\Delta,\epsilon)$ and assume $T(\mathbf x)<T(\mathbf y)$.
Set $\mathbf x':=\mathbf\Phi(\mathbf x,T(\mathbf x))$ and $\mathbf y':=\mathbf\Phi(\mathbf y,T(\mathbf x))$.
Then we have $|\mathbf x'\mathbf y'|\le C|\mathbf x\mathbf y|$ and thus $|\Delta\mathbf y'|\le C|\mathbf x\mathbf y|$.
Since $\mathbf\Phi$ decreases the distance to $\Delta$ with velocity almost $1$ (Lemma \ref{lem:reg}), we obtain $T(\mathbf y')\le2C|\mathbf x\mathbf y|$.
Therefore
\[T(\mathbf y)\le T(\mathbf x)+T(\mathbf y')\le T(\mathbf x)+2C|\mathbf x\mathbf y|.\]

Now we define $\mathbf\Psi:\mathcal R(\delta,\ell)^2\cap B(\Delta,\epsilon)\times[0,\epsilon]\to B(\Delta,\epsilon)$ by
\[\mathbf\Psi(\mathbf x,t):=\mathbf\Phi(\mathbf x,(t/\epsilon)T(\mathbf x)).\]
It is easy to check that $\mathbf\Psi$ satisfies the conclusion.
\end{proof}

Let us write $\mathbf\Psi=(\Psi_1,\Psi_2)\in M\times M$.
For $(x,y)\in\mathcal R(\delta,\ell)^2\cap B(\Delta,\epsilon)$, by joining the curves $\Psi_1(x,y,2t)$ and $\Psi_2(x,y,\epsilon-2t)$, we obtain

\begin{cor}\label{cor:diag}
There is a $C$-Lipschitz map $\Psi:\mathcal R(\delta,\ell)^2\cap B(\Delta,\epsilon)\times[0,\epsilon]\to M$, where $C$ is independent of any a priori constant, such that
\[\Psi(x,y,0)=x,\quad\Psi(x,y,\epsilon)=y,\quad\Psi(x,x,t)=x\]
for any $(x,y)\in\mathcal R(\delta,\ell)^2\cap B(\Delta,\epsilon)$ and $0\le t\le\epsilon$.
\end{cor}

\begin{proof}[Proof of Theorem \ref{thm:isom}]
Assume that $M$ and $M'$ are $(C^{-1}\epsilon)$-close in the Gromov-Hausdorff distance, where $\epsilon<C^{-1}\delta\ell$.
By Theorem \ref{thm:main} there exists a $(C,\epsilon)$-Lipschitz homotopy approximation $f:M\to M'$ with inverse $f':M'\to M$.
Let $g:\mathcal R_M(10\delta,\ell/10)\to M'$ be a $\tau(\delta)$-almost isometric open embedding as in Theorem \ref{thm:reg}.
Note that by construction $f$ and $g$ are $\epsilon$-close on their common domain and also that the image of $g$ contains $\mathcal R_{M'}(9\delta,\ell/9)$ (the latter follows from the openness of $g$).

We first glue these two maps by using Corollary \ref{cor:diag} (one can also use the center of mass technique as in \cite{MY:lip}).
Let $\Psi':\mathcal R_{M'}(10\delta,\ell/10)^2\cap B(\Delta,C\epsilon)\times[0,\epsilon]\to M'$ be a $C$-Lipschitz map as in Corollary \ref{cor:diag} (we are rescaling the time interval as usual).
We define $h:M\to M'$ by
\[
h(x):=
\begin{cases}
\Psi'(g(x),f(x),\min\{|\mathcal R_M(2\delta,\ell/2),x|,\epsilon\}) & x\in\mathcal R_M(3\delta,\ell/3)\\
\hfil f(x) &  \text{otherwise}.
\end{cases}
\]
Then $h$ is a $C$-Lipschitz $C\epsilon$-approximation which coincides with $g$ on $\mathcal R_M(2\delta,\ell/2)$ and with $f$ outside $\mathcal R_M(3\delta,\ell/3)$.
Similarly we define $h':M'\to M$ by gluing $f'$ and $g^{-1}$ by using $\Psi:\mathcal R_{M}(10\delta,\ell/10)^2\cap B(\Delta,C\epsilon)\times[0,\epsilon]\to M$ so that $h'=g^{-1}$ on $\mathcal R_{M'}(2\delta,\ell/2)$ and $h'=f'$ outside $\mathcal R_{M'}(3\delta,\ell/3)$.
Then we have
\[
h'\circ h(x)=
\begin{cases}
\hfil x & x\in\mathcal R_M(\delta,\ell) \\
f'\circ f(x) & x\notin\mathcal R_M(4\delta,\ell/4).
\end{cases}
\]

Let $F$ be a $(C,\epsilon)$-Lipschitz homotopy between $f'\circ f$ and $1_M$.
Set $G(x,t):=\Psi(h'\circ h(x),x,t)$, which is a $(C,\epsilon)$-Lipschitz homotopy between $h'\circ h$ and $1_M$ on $\mathcal R_M(9\delta,\ell/9)$, fixing $\mathcal R_M(\delta,\ell)$.
Since $h'\circ h=f'\circ f$ outside $\mathcal R_M(4\delta,\ell/4)$, one can glue these two homotopies by using $\Psi$, that is, we define $H:M\times[0,\epsilon]\to M$ by
\[H(x,t)
:=
\begin{cases}
\Psi(G(x,t),F(x,t),\min\{|\mathcal R_M(4\delta,\ell/4),x|,\epsilon\}) & x\in\mathcal R_M(5\delta,\ell/5)\\
\hfil F(x,t) & \text{otherwise}.
\end{cases}
\]
Then $H$ is a $(C,\epsilon)$-Lipschitz homotopy between $h'\circ h$ and $1_M$ which fixes $\mathcal R_M(\delta,\ell)$.
Similarly one can construct a $(C,\epsilon)$-Lipschitz homotopy between $h\circ h'$ and $1_{M'}$.
This completes the proof.
\end{proof}

\section{Proof of Theorem \ref{thm:ext}}\label{sec:ext}

In this section we give two proofs of Theorem \ref{thm:ext}.
The first one is a direct corollary of the proofs of Theorem \ref{thm:main} and the second one provides a more general way of deforming a Lipschitz homotopy approximation.

We first define extremal subsets.
Here we state the definition in \cite{Pet:sc} by gradient flows, which is more suitable for our application than the original one in \cite{PP:ext}.

\begin{dfn}\label{dfn:ext}
A subset $E$ of an Alexandrov space $M$ is said to be \textit{extremal} if the gradient curve of any semiconcave function starting at a point of $E$ remains in $E$.
\end{dfn}

Note that any extremal subset is closed.
A typical example is the boundary of an Alexandrov space.
See \cite[\S4]{Pet:sc} or \cite{PP:ext} for other examples and properties of extremal subsets.

\subsection{First Proof}

We prove that the Lipschitz homotopies constructed in Sections \ref{sec:main1} and \ref{sec:main2} automatically preserve extremal subsets (a slight modification is needed for the latter one).
For each statement in Sections \ref{sec:main1} and \ref{sec:main2}, we will give a complement concerning extremal subsets.
The key fact is that the $(C,\epsilon)$-Lipschitz contraction of Theorem \ref{thm:conv} was given by gradient flows, and hence it automatically preserves extremal subsets.
Note that the contents of this subsection is independent of Lipschitz continuity.

\begin{conv}
In the following complements, we always assume that $M$ and $M'$ are Alexandrov spaces in $\mathcal A(n,D,v)$ (not general metric spaces as in Sections \ref{sec:main1} and \ref{sec:main2}).
We also assume that $(C,\epsilon)$-Lipschitz contractible domains and $(C,\epsilon)$-good coverings are the ones obtained by the proofs of Theorems \ref{thm:conv} and \ref{thm:good} such that the contractions are defined by gradient flows.
\end{conv}

\subsubsection{Complements to Section \ref{sec:main1}}

Let $M\in\mathcal A(n,D,v)$ and let $\mathcal U=\{U_j\}_j$ be a $(C,\epsilon)$-good covering of $M$ given by Theorem \ref{thm:good}.
We use the same notation as in Section \ref{sec:main1}: $\mathcal N$ and $|\mathcal N|$ denote the nerve of $\mathcal U$ and its $\epsilon$-geometric realization, respectively.
Furthermore, for $\sigma\in\mathcal N$, we denote by $U_\sigma$ the corresponding intersection of $U_j$ and by $p_\sigma$ its center.

Let $E\subset M$ be an extremal subset.
We define a subcomplex $\mathcal N(E)$ of $\mathcal N$ as a subset consisting of $\sigma\in\mathcal N$ such that $U_\sigma$ intersects $E$.
We identify $|\mathcal N(E)|$ with a subset of $|\mathcal N|$.

Since the Lipschitz contraction of Theorem \ref{thm:good}(3) is given by gradient flows, we have

\begin{lem}\label{lem:cen}
If $\sigma\in\mathcal N(E)$, then $p_\sigma\in E$.
\end{lem}

Using this lemma, we prove

\begin{comp}[to Theorem \ref{thm:nerve}]\label{comp:nerve}
$(M,E)$ is $(\tilde C,\epsilon)$-Lipschitz homotopy equivalent to $(|\mathcal N|,|\mathcal N(E)|)$.
\end{comp}

\begin{proof}
Let $\Theta:M\to|\mathcal N|$ and $\zeta:|\mathcal N|\to M$ be the $(\tilde C,\epsilon)$-Lipschitz homotopy equivalences constructed in the proof of Theorem \ref{thm:nerve}.
We prove that they and their homotopies respect $E$ and $|\mathcal N(E)|$.

It is clear that $\Theta(E)\subset |\mathcal N(E)|$.
To show $\zeta(|\mathcal N(E)|)\subset E$, let us observe how the map $\zeta$ is constructed.
Recall that $\zeta=q\circ\bar\Phi\circ\Psi$ (see \eqref{eq:zeta}), where $q$ is a projection, $\Psi$ is an embedding, and $\bar\Phi=\Phi(\cdot,\epsilon)$ is a retraction which is a composition of $\bar\Phi_\sigma=\Phi_\sigma(\cdot,\epsilon)$ as in Subclaim \ref{subclm:Phi}.
In what follows, we will focus only on the ``$M$-coordinate'' in $\mathcal M$.
More precisely, if $y=(\theta,[x,t])\in\mathcal M\subset |\mathcal N|\times K(M)$ and if $t\neq\epsilon$, then we call $x$ the \textit{$M$-coordinate} of $y$.

Let $\theta\in|\mathcal N|$ and set $\sigma:=\supp\theta$.
By definition $\Psi(\theta)=(\theta,v_M)$, where $v_M$ is the vertex of $K(M)$.
Depending on $\theta$, there are two possibilities for its $\bar\Phi_\sigma$-image:
\begin{enumerate}
\item the $M$-coordinate of $\bar\Phi_\sigma(\theta,v_M)$ is $p_\sigma$; or
\item $\bar\Phi_\sigma(\theta,v_M)=(\eta,v_M)$ for some $\eta\in\partial\sigma$
\end{enumerate}
(see \eqref{eq:varphi} and \eqref{eq:r}; in fact there is no natural candidate for the $M$-coordinate other than $p_\sigma$ since $(\theta,v_M)$ has no $M$-coordinate).
Let us next consider $\bar\Phi_\tau\circ\bar\Phi_\sigma(\theta,v_M)$, where $\tau$ is a codimension $1$ face of $\sigma$ for which this composition makes sense.
In the case (1), we see that $p_\sigma$ is moved by the Lipschitz contraction of $U_\tau$ to the $M$-coordinate of $\bar\Phi_\tau\circ\bar\Phi_\sigma(\theta,v_M)$.
In the case (2), we can repeat the same argument.

By repeating this observation, we see that in the $M$-coordinate, $\theta$ is first mapped to some $p_{\sigma'}$, where $\sigma'$ is a face of $\sigma$, and then moved by the Lipschitz contractions of $U_\tau$ to reach $\zeta(\theta)$, where $\tau$ are faces of $\sigma'$.
In particular, if $\theta\in|\mathcal N(E)|$ then $p_{\sigma'}\in E$ by Lemma \ref{lem:cen}.
Since $E$ is closed under gradient flows, this implies $\zeta(\theta)\in E$.

Next we show that the homotopy between $\zeta\circ\Theta$ and $1_M$ respects $E$.
Observe that this homotopy is essentially given by Claim \ref{clm:Psi}.
More precisely, for any $x\in M$, the homotopy curve is defined by
\[s\mapsto q\circ\bar\Phi(\Theta(x),[x,s]).\]
The above argument shows that this is a broken gradient curve from $x$ to $\zeta\circ\Theta(x)$.
Therefore if $x\in E$ it is entirely contained in $E$.

Finally we show that the homotopy between $\Theta\circ\zeta$ and $1_{|\mathcal N|}$ respects $|\mathcal N(E)|$.
Observe that this homotopy is essentially given by Claims \ref{clm:tau} and \ref{clm:Phi}.
For any $\theta\in|\mathcal N|$, we see that the homotopy curve between $\theta$ and $\Theta\circ\zeta(\theta)$ is a broken line through $p\circ\bar\Phi(\theta,v_M)$, where $p$ is the projection to $|\mathcal N|$.
More precisely, the homotopy of Claim \ref{clm:tau} gives a line segment between $p\circ\bar\Phi(\theta,v_M)$ and $\Theta\circ\zeta(\theta)$ defined by
\[s\mapsto p\circ H(\bar\Phi(\theta,v_M),s).\]
On the other hand, the homotopy of Claim \ref{clm:Phi} gives a broken line joining $\theta$ and $p\circ\bar\Phi(\theta,v_M)$ defined by
\[s\mapsto p\circ\Phi((\theta,v_M),s),\]
which lies in $\supp\theta$.
Therefore if $\theta\in|\mathcal N(E)|$ the homotopy curve is entirely contained in $|\mathcal N(E)|$.
\end{proof}

To prove Theorem \ref{thm:ext}, we also need the following stability.

\begin{comp}[to Theorem \ref{thm:good}(4)]\label{comp:good}
Let $E$ and $E'$ be extremal subsets of $M$ and $M'$, respectively, that are $(C^{-1}\epsilon)$-close to each other.
Then the subcomplexes of the nerves of $\mathcal U$ and $\mathcal U'$ defined by $E$ and $E'$ are the same.
\end{comp}

\begin{proof}
Let $U$ be an intersection of elements of $\mathcal U$ and let $U'$ denote its lift to $M'$.
We show that $U$ intersects $E$ if and only if  $U'$ intersects $E'$.
Let $p$ and $p'$ denote the peaks of $U$ and $U'$ (different from the centers of the contractions).
By Claim \ref{clm:good}, $U$ and $U'$ contain the $(C_1^{-1}\epsilon)$-balls around $p$ and $p'$, respectively.
Furthermore, $p$ is $(C_2^{-1}\epsilon)$-close to $p'$, where $C_1\ll C_2\ll C$; otherwise by the gradient estimate (Lemma \ref{lem:grad}) we get a contradiction.

Suppose $U$ intersects $E$.
It also follows from the gradient estimate and the extremality of $E$ that $p\in E$.
Therefore there exists a point of $E'$ that is $(2C_2^{-1}\epsilon)$-close to $p'$.
Since $U'$ contains the $(C_1^{-1}\epsilon)$-ball around $p'$, this shows that $U'$ intersects $E'$.
The converse is proved in exactly the same way.
\end{proof}

Theorem \ref{thm:ext} immediately follows from Complements \ref{comp:nerve} and \ref{comp:good} in the same way as Theorem \ref{thm:main} in Section \ref{sec:main1}.

\subsubsection{Complements to Section \ref{sec:main2}}

Let $M\in\mathcal A(n,D,v)$ and $E$ an extremal subset.
By the proofs and the fact that $(C,\epsilon)$-Lipschitz contractions preserve extremal subsets, we have

\begin{comp}[to Proposition \ref{prop:simp}]\label{comp:simp}
Let $\sigma\in K$ be a simplex such that $f(v)\in E$ for any vertex $v\in\sigma$ and $f(\sigma\cap|L|)\subset E$.
Then the extension $\tilde f$ satisfies $\tilde f(\sigma)\subset E$.
\end{comp}

\begin{comp}[to Corollary \ref{cor:simp}]\label{comp:simp'}
Let $L$ be a subcomplex of $K$ such that $f_i(|L|)\subset E$ for $i=0,\epsilon$.
Then the obtained homotopy respects $|L|$ and $E$, i.e., the images of $|L|$ are always contained in $E$.
\end{comp}

As in Subsection \ref{sec:dom}, let $\mathcal N$ be the nerve of a $2\epsilon$-covering of $M$ centered at $\epsilon/2$-discrete points and $|\mathcal N|$ its $\epsilon$-geometric realization.
As before, we define a subcomplex $\mathcal N(E)$ of $\mathcal N$ as a subset consisting of simplices such that the corresponding intersections of $2\epsilon$-balls intersect $E$.

The following is proved in the same way as Complement \ref{comp:nerve}.

\begin{comp}[to Proposition \ref{prop:dom}]\label{comp:dom}
$(M,E)$ is $(\tilde C,\epsilon)$-Lipschitz dominated by $(|\mathcal N|,|\mathcal N(E)|)$, i.e., the maps $\Theta$ and $\zeta$, and the homotopy between $\zeta\circ\Theta$ and $1_M$ respect $E$ and $|\mathcal N(E)|$.
\end{comp}

By Complements \ref{comp:simp'} and \ref{comp:dom} we have

\begin{comp}[to Proposition \ref{prop:hom}]\label{comp:hom}
Let $E$ and $E'$ be extremal subsets in $M$ and $M'$, respectively, such that $f_i(E)\subset E'$ for $i=0,\epsilon$.
Then the obtained homotopy respects $E$ and $E'$.
\end{comp}

Theorem \ref{thm:ext} follows from Complements \ref{comp:simp}, \ref{comp:dom} and \ref{comp:hom} in almost the same way as Theorem \ref{thm:main} in Section \ref{sec:main2}.
The only modification is as follows.
Let $M,M'\in\mathcal A(n,D,v)$ and $E\subset M$, $E'\subset M'$ extremal subsets that are sufficiently close in the Gromov-Hausdorff distance.
When we construct a map $F=f\circ\Theta:M\to M'$ as in Section \ref{sec:main2}, we choose $\{p_j'\}_j\subset M'$ so that $p_j'\in E'$ whenever $B(p_j,2\epsilon)$ intersects $E$.
Then Complements \ref{comp:simp} and \ref{comp:dom} imply that $F(E)\subset E'$.
Similarly we get a map $F':M'\to M$ satisfying $F'(E')\subset E$.
The rest follows from Complement \ref{comp:hom}.

\subsection{Second Proof}

We prove that any Lipschitz homotopy approximation between Alexandrov spaces can be deformed to preserve close extremal subsets.

The proof is based on the following theorem.
Fix $n$, $D$, and $v$.

\begin{thm}[{\cite[1.5]{F:reg}}]\label{thm:ndr}
There is $C=C(n,D,v)>0$ satisfying the following:
Let $M\in\mathcal A(n,D,v)$ and $E\subset M$ an extremal subset.
Then for any $0<\epsilon<C^{-1}$, the $\epsilon$-neighborhood $B(E,\epsilon)$ of $E$ admits a $(C,\epsilon)$-Lipschitz strong deformation retraction to $E$.
\end{thm}

Although the original statement in \cite[1.5]{F:reg} does not assert the existence of such a uniform constant, one can easily check it by the proof together with the following lemma.

\begin{lem}\label{lem:ndr}
Let $M\in\mathcal A(n,D,v)$ and $E\subset M$ a proper extremal subset.
Then there exists $q\in M$ such that $|qE|>c(n,D,v)$, where $c(n,D,v)$ is a positive constant depending only on $n$, $D$, and $v$.
\end{lem}

\begin{proof}
Any limit of proper extremal subsets under a noncollapsing convergence of Alexandrov spaces is a proper extremal subset (\cite[9.13]{K:stab}).
Hence the statement follows from the compactness of $\mathcal A(n,D,v)$.
\end{proof}

\begin{proof}[Proof of Theorem \ref{thm:ndr}]
We explain how to modify the original proof.
In what follows we denote by $C$ and $c$ various large and small positive constants depending only on $n$, $D$, and $v$.
Recall that the Lipschitz deformation retraction of \cite[1.5]{F:reg} was obtained by modifying the gradient flow $\Phi$ of a function
\[h=\frac1N\sum_{\alpha=1}^N|q_\alpha\cdot|,\]
where $\{q_\alpha\}_{\alpha=1}^N$ is a $\nu$-discrete set in $M\setminus E$ such that $N\ge c/\nu^n$ for sufficiently small $\nu>0$.
By Lemma \ref{lem:ndr}, $\{q_\alpha\}_{\alpha=1}^N$ can be chosen so that $\rho:=\min_\alpha|q_\alpha E|>c$.
Then in the $\rho/2$-neighborhood of $E$ the concavity of $h$ depends only on $n$, $D$, and $v$.
Hence the gradient flow $\Phi$ is uniformly $C$-Lipschitz in this neighborhood (as long as the time is uniformly bounded above; Proposition \ref{prop:lip}).

In \cite[A.1]{F:reg} we showed that $d\dist_E(\nabla h)<-c$ on $B(E,\epsilon)\setminus E$, where $\epsilon\ll\rho$.
Since the proof relies only on volume comparison, one can check that the upper bound of $\epsilon$ depends only on $n$, $D$, and $v$.

The rest of the proof is the same as that of Proposition \ref{prop:diag}.
Let $T(x)$ denote the minimum time such that the gradient flow $\Phi$ starting at $x\in B(E,\epsilon)$ reaches $E$.
From the above, $T(x)$ is uniformly bounded above by $C\epsilon$.

In \cite[A.2]{F:reg} we showed that the function $x\mapsto T(x)$ is Lipschitz.
Since the proof relies only on \cite[A.1]{F:reg} and the Lipschitz continuity of the gradient flow, one can check that it is uniformly $C$-Lipschitz.

The desired deformation retraction $\Psi:B(E,\epsilon)\times[0,\epsilon]\to B(E,\epsilon)$ is given by
\[\Psi(x,t):=\Phi\left(x,(t/\epsilon)T(x)\right).\]
A direct calculation shows that $\Psi$ is also $C$-Lipschitz.
\end{proof}

\begin{rem}\label{rem:ndr}
As can be seen in the proof, the flow $\Psi$ monotonically decreases the distance to $E$.
\end{rem}

\begin{proof}[Proof of Theorem \ref{thm:ext}]
As before, we denote by $C$ various large positive constants depending only on $n$, $D$, and $v$.
Let $M,M'\in\mathcal A(n,D,v)$ and $E\subset M$, $E'\subset M'$ extremal subsets that are $(C^{-1}\epsilon)$-close in the Gromov-Hausdorff distance, where $0<\epsilon<C^{-1}$.
By Theorem \ref{thm:main}, there exist $(C,\epsilon)$-Lipschitz homotopy approximations $f:M\to M'$ and $f':M'\to M$ with $(C,\epsilon)$-Lipschitz homotopies $F:M\times[0,\epsilon]\to M$ and $F':M'\times[0,\epsilon]\to M'$ (not necessarily the ones given in the proof of Theorem \ref{thm:main}).
Taking the Gromov-Hausdorff distance smaller (by multiplying by a constant) if necessary, we may assume that
\begin{gather*}
f(E)\subset B(E',\epsilon),\quad f'(E')\subset B(E,\epsilon),\\
F(E,t)\subset B(E,\epsilon),\quad F'(E',t)\subset B(E',\epsilon)
\end{gather*}
for all $0\le t\le\epsilon$.
We deform these maps to respect $E$ and $E'$.

Let $\Psi$ and $\Psi'$ be the $(C,\epsilon)$-Lipschitz strong deformation retractions of $B(E,10\epsilon)$ and $B(E',10\epsilon)$ to $E$ and $E'$ provided by Theorem \ref{thm:ndr}, respectively.
We extend $\Psi$ to $\tilde\Psi:M\times[0,\epsilon]\to M$ as follows:
\[
\tilde\Psi(x,t):=
\begin{cases}
\Psi(x,\rho(|Ex|)t) & x\in B(E,10\epsilon)\\
\hfil x & \text{otherwise},
\end{cases}
\]
where $\rho:[0,\infty)\to[0,1]$ is a $(5\epsilon)^{-1}$-Lipschitz function such that $\rho\equiv1$ on $[0,5\epsilon]$ and $\rho\equiv0$ on $[10\epsilon,\infty)$.
It is straightforward to check that $\tilde \Psi$ is $C$-Lipschitz independent of $\epsilon$ and $\tilde\Psi(\cdot,0)=1_M$.
We also define $\tilde\Psi':M'\times[0,\epsilon]\to M'$ in the same manner.
Let $\tilde\Psi_t$ and $\tilde\Psi_t'$ denote $\tilde\Psi(\cdot,t)$ and $\tilde\Psi'(\cdot,t)$, respectively.

Now we define new maps $\tilde f:M\to M'$ and $\tilde f':M'\to M$ by
\[\tilde f:=\tilde \Psi_\epsilon'\circ f,\quad\tilde  f':=\tilde\Psi_\epsilon\circ f'.\]
Then they are $C$-Lipschitz $10\epsilon$-approximations satisfying $\tilde f(E)\subset E'$ and $\tilde f'(E')\subset E$.
Moreover,
\[\tilde f'\circ\tilde f=\tilde\Psi_\epsilon\circ f'\circ\tilde \Psi'_\epsilon\circ f\sim\tilde\Psi_\epsilon\circ f'\circ f\sim\tilde\Psi_\epsilon\sim1_M,\]
where $\sim$ denote $(C,\epsilon)$-Lipschitz homotopies.
It is easy to see that these homotopies respect $E$ (note $\tilde\Psi_t'\circ f(E)\subset B(E',\epsilon)$ by Remark \ref{rem:ndr}).
Similarly $\tilde f\circ\tilde f'$ is $(C,\epsilon)$-Lipschitz homotopic to $1_{M'}$ while preserving $E'$, which completes the proof.
\end{proof}

\begin{thebibliography}{99}

\bibitem{BGP}
Yu. Burago, M. Gromov, and G. Perel'man, A. D. Aleksandrov spaces with curvatures bounded below, Uspekhi Mat. Nauk 47 (1992), no. 2(284), 3--51, 222; translation in Russian Math. Surveys 47 (1992), no. 2, 1--58.

\bibitem{F:reg}
T. Fujioka, Regular points of extremal subsets in Alexandrov spaces, J. Math Soc. Japan 74 (2022), no. 4, 1245--1268.

\bibitem{K:stab}
V. Kapovitch, Perelman's stability theorem, Metric and comparison geometry, 103--136, Surv. Differ. Geom., 11, Int. Press, Somerville, MA, 2007.

\bibitem{MY:loc}
A. Mitsuishi and T. Yamaguchi, Locally Lipschitz contractibility of Alexandrov spaces and its applications, Pacific J. Math. 270 (2014), no. 2, 393--421.

\bibitem{MY:stab}
A. Mitsuishi and T. Yamaguchi, Stability of strongly Lipschitz contractible balls in Alexandrov spaces, Math. Z. 277 (2014), 995--1009. 

\bibitem{MY:good}
A. Mitsuishi and T. Yamaguchi, Good coverings of Alexandrov spaces, Trans. Amer. Math. Soc. 372 (2019), 8107--8130.

\bibitem{MY:lip}
A. Mitsuishi and T. Yamaguchi, Lipschitz homotopy convergence of Alexandrov spaces, J. Geom. Anal. 29 (2019), 2217--2241.

\bibitem{Per:alex}
G. Perelman, Alexandrov's spaces with curvatures bounded from below II, preprint, 1991.

\bibitem{Per:mor}
G. Ya. Perel'man, Elements of Morse theory on Aleksandrov spaces, Algebra i Analiz 5 (1993), no. 1, 232--241; translation in St. Petersburg Math. J. 5 (1994), no. 1, 205--213.

\bibitem{PP:ext}
G. Ya. Perel'man and A. M. Petrunin, Extremal subsets in Aleksandrov spaces and the generalized Liberman theorem, Algebra i Analiz 5 (1993), no. 1, 242--256; translation in St. Petersburg Math. J. 5 (1994), no. 1, 215--227.

\bibitem{Peter}
P. Petersen V, A finiteness theorem for metric spaces, J. Differential Geom. 31 (1990), no. 2, 387--395.

\bibitem{Pet:sc}
A. Petrunin, Semiconcave functions in Alexandrov's geometry, Metric and comparison geometry, Metric and comparison geometry, 137--201, Surv. Differ. Geom., 11, Int. Press, Somerville, MA, 2007.

\bibitem{Sha}
A. R. Shastri, Basic Algebraic Topology, CRC Press, Taylor \& Francis Group, A Chapman \& Hall Book, New York, 2014.

\bibitem{WSS}
Y. Wang, X. Su and H. Sun, A new proof of almost isometry theorem in Alexandrov geometry with curvature bounded below, Asian J. Math. 17 (2013), no. 4, 715--728.

\bibitem{Y:conv}
T. Yamaguchi, A convergence theorem in the geometry of Alexandrov spaces, Actes de la table ronde de g\'eom\'etrie diff\'erentielle, 601--642, S\'emin. Congr., 1, Soc. Math. France, Paris, 1996.

\bibitem{Y:ess}
T. Yamaguchi, Collapsing and essential coverings, preprint, arXiv:1205.0441, 2012.

\end{thebibliography}

\end{document}
